\documentclass[UTF8,fontset=fandol,openany]{ctexbook}
% ctexbook defaults are retained; only page geometry and technical support are added.
\usepackage[a4paper,inner=18mm,outer=16mm,top=18mm,bottom=19mm,headheight=15pt]{geometry}
\usepackage{amsmath,amssymb,amsthm,mathtools}
\usepackage{graphicx,tikz,adjustbox}
\usetikzlibrary{arrows.meta,positioning,calc,fit,shapes.geometric,decorations.pathreplacing,decorations.pathmorphing}
\usepackage{array,longtable,booktabs,multirow,tabularx}
\usepackage[normalem]{ulem}
\usepackage{xeCJKfntef}
\xeCJKDeclareCharClass{CJK}{"2460 -> "2473}
\usepackage{enumitem}
\usepackage{hyperref}
\hypersetup{unicode=true,hidelinks,bookmarksnumbered=true,pdfcreator={XeLaTeX / ctexbook}}
\usepackage{bookmark}
\setcounter{tocdepth}{1}
\setcounter{secnumdepth}{2}
\let\cleardoublepage\clearpage
% Large nonbreaking TikZ figures may move to the next page. Keep ordinary
% paragraph spacing instead of stretching the preceding page vertically.
\raggedbottom
\setlength{\emergencystretch}{3em}
\newcommand{\SourceTag}[1]{\textbf{【#1】}}
\newcommand{\SourceFootnote}[1]{\par\smallskip\noindent{\footnotesize #1}\par\smallskip}
\newcommand{\InsertDiagram}[1]{\par\smallskip\begin{center}\begin{adjustbox}{max width=\linewidth,max totalheight=.78\textheight}\input{#1}\end{adjustbox}\end{center}\smallskip}
\newcommand{\OriginalPanel}[1]{\par\smallskip\noindent\textbf{#1}\par}
\newcommand{\BookEditionDate}{2026 年 11 月 20 日}
\newcommand{\BookCover}[4]{%
  \begin{titlepage}\centering
  \vspace*{20mm}
  {\zihao{2}\bfseries #1\par}
  \vspace{10mm}
  {\zihao{-0}\bfseries 刑法\par}
  \vspace{8mm}
  {\zihao{2}#2\par}
  \vspace{12mm}
  {\zihao{4}#3\par}
  \vfill
  {\zihao{-4}#4\par}
  \vspace{9mm}
  {\zihao{-4}\BookEditionDate\par}
  \vspace*{14mm}
  \end{titlepage}}

\begin{document}
\setbox0=\vbox{\hsize=2.640cm\parindent=0pt\linespread{1}\fontsize{10}{13.50}\selectfont\bfseries\centering\noindent 误解的类型\par}
\typeout{TCELL:p303_2:0:\the\ht0:\the\dp0}
\setbox0=\vbox{\hsize=4.140cm\parindent=0pt\linespread{1}\fontsize{10}{13.50}\selectfont\bfseries\centering\noindent 具体情形\par}
\typeout{TCELL:p303_2:1:\the\ht0:\the\dp0}
\setbox0=\vbox{\hsize=8.440cm\parindent=0pt\linespread{1}\fontsize{10}{13.50}\selectfont\bfseries\centering\noindent 范例\par}
\typeout{TCELL:p303_2:2:\the\ht0:\the\dp0}
\setbox0=\vbox{\hsize=2.640cm\parindent=0pt\linespread{1}\fontsize{10}{13.50}\selectfont\raggedright\noindent 误解一：只要客观上具备数额，就定罪量刑\par}
\typeout{TCELL:p303_2:3:\the\ht0:\the\dp0}
\setbox0=\vbox{\hsize=4.140cm\parindent=0pt\linespread{1}\fontsize{10}{13.50}\selectfont\raggedright\noindent （1）只要客观上达到“数额较大”，就定罪\par}
\typeout{TCELL:p303_2:4:\the\ht0:\the\dp0}
\setbox0=\vbox{\hsize=8.440cm\parindent=0pt\linespread{1}\fontsize{10}{13.50}\selectfont\raggedright\noindent 农民甲以为自己盗窃的是普通葡萄，实际是高科技试验品，价值40万元。有意见认为，既然该葡萄是数额特别巨大的财物（当然也算数额较大的财物），因此甲构成盗窃罪。然而，由于甲没有认识到该葡萄是数额较大的财物，因此不构成盗窃罪\par}
\typeout{TCELL:p303_2:5:\the\ht0:\the\dp0}
\setbox0=\vbox{\hsize=4.140cm\parindent=0pt\linespread{1}\fontsize{10}{13.50}\selectfont\raggedright\noindent （2）只要客观上达到“数额（特别）巨大”，就适用“数额（特别）巨大”的法定刑\par}
\typeout{TCELL:p303_2:6:\the\ht0:\the\dp0}
\setbox0=\vbox{\hsize=8.440cm\parindent=0pt\linespread{1}\fontsize{10}{13.50}\selectfont\raggedright\noindent 客人王某带一盆兰花入住五星级酒店时，服务员甲认为这盆兰花应该值几千元，便盗窃到手。物价部门鉴定，该兰花属于稀世珍品，价值100万元。有意见认为，既然该兰花是数额特别巨大的财物，因此，对甲应适用“数额特别巨大”档次的法定刑。\\{}然而，甲虽然认识到该兰花是数额较大的财物，但是没有认识到是数额特别巨大的财物，因此，只能适用“数额较大”档次的法定刑，不能适用“数额特别巨大”档次的法定刑\par}
\typeout{TCELL:p303_2:7:\the\ht0:\the\dp0}
\setbox0=\vbox{\hsize=2.640cm\parindent=0pt\linespread{1}\fontsize{10}{13.50}\selectfont\bfseries\centering\noindent 误解的类型\par}
\typeout{TCELL:p304_1:0:\the\ht0:\the\dp0}
\setbox0=\vbox{\hsize=4.140cm\parindent=0pt\linespread{1}\fontsize{10}{13.50}\selectfont\bfseries\centering\noindent 具体情形\par}
\typeout{TCELL:p304_1:1:\the\ht0:\the\dp0}
\setbox0=\vbox{\hsize=8.440cm\parindent=0pt\linespread{1}\fontsize{10}{13.50}\selectfont\bfseries\centering\noindent 范例\par}
\typeout{TCELL:p304_1:2:\the\ht0:\the\dp0}
\setbox0=\vbox{\hsize=2.640cm\parindent=0pt\linespread{1}\fontsize{10}{13.50}\selectfont\raggedright\noindent 误解二：只要客观上不具备数额，就不定罪量刑\par}
\typeout{TCELL:p304_1:3:\the\ht0:\the\dp0}
\setbox0=\vbox{\hsize=4.140cm\parindent=0pt\linespread{1}\fontsize{10}{13.50}\selectfont\raggedright\noindent （1）只要客观上未达到“数额较大”，就不定罪。然而，盗窃罪的成立条件是，主观上有盗窃数额较大财物的故意，客观上有盗窃到数额较大财物的可能性、危险性。既遂条件是，盗窃到数额较大的财物\par}
\typeout{TCELL:p304_1:4:\the\ht0:\the\dp0}
\setbox0=\vbox{\hsize=8.440cm\parindent=0pt\linespread{1}\fontsize{10}{13.50}\selectfont\raggedright\noindent 甲在商场盗窃到女士乙的提包，发现里面只有100元（不考虑提包本身的价值）。有意见认为，由于该财物未达到数额较大，因此甲不构成盗窃罪。然而，甲应成立盗窃罪，因为甲有盗窃到数额较大财物的可能性、危险性。甲构成犯罪未遂\par}
\typeout{TCELL:p304_1:5:\the\ht0:\the\dp0}
\setbox0=\vbox{\hsize=4.140cm\parindent=0pt\linespread{1}\fontsize{10}{13.50}\selectfont\raggedright\noindent （2）只要客观上未达到“数额（特别）巨大”，就不适用“数额（特别）巨大”的法定刑。然而，多数说认为，只要主观上有盗窃数额巨大财物的故意，客观上有盗窃到数额巨大财物的可能性，也即对该财物具有危险性，那么就应适用“数额巨大”的法定刑；若未盗窃到，再适用未遂规定\par}
\typeout{TCELL:p304_1:6:\the\ht0:\the\dp0}
\setbox0=\vbox{\hsize=8.440cm\parindent=0pt\linespread{1}\fontsize{10}{13.50}\selectfont\raggedright\noindent 甲欲盗窃某珠宝店3万元的珠宝，手伸进玻璃柜（此时已经属于盗窃着手），左边是3万元的珠宝，右边是3千元的珠宝。甲认错对象，拿走3千元的珠宝。（1）多数说（司法解释）认为，甲应构成盗窃“数额巨大财物”的未遂，①同时，甲构成盗窃“数额较大财物”的既遂。甲的一个行为符合两个法定刑条件，想象竞合，择一重刑论处。理由是，盗窃“数额巨大财物”这种法定刑升格条件属于加重型构成要件，存在既遂与未遂问题。（2）少数说认为，由于甲未盗窃到数额巨大的财物，因此不能适用“数额巨大”档次的法定刑。理由是，“数额巨大”只有客观具备与否问题，没有既遂与未遂问题，属于单纯的量刑规则。对甲只能适用盗窃“数额较大财物”的既遂\par}
\typeout{TCELL:p304_1:7:\the\ht0:\the\dp0}
\setbox0=\vbox{\hsize=2.440cm\parindent=0pt\linespread{1}\fontsize{10}{13.50}\selectfont\bfseries\centering\noindent 非法占有目的\par}
\typeout{TCELL:p307_1:0:\the\ht0:\the\dp0}
\setbox0=\vbox{\hsize=1.540cm\parindent=0pt\linespread{1}\fontsize{10}{13.50}\selectfont\bfseries\centering\noindent 行为对象\par}
\typeout{TCELL:p307_1:1:\the\ht0:\the\dp0}
\setbox0=\vbox{\hsize=1.540cm\parindent=0pt\linespread{1}\fontsize{10}{13.50}\selectfont\bfseries\centering\noindent 行为方式\par}
\typeout{TCELL:p307_1:2:\the\ht0:\the\dp0}
\setbox0=\vbox{\hsize=2.290cm\parindent=0pt\linespread{1}\fontsize{10}{13.50}\selectfont\bfseries\centering\noindent 结论\par}
\typeout{TCELL:p307_1:3:\the\ht0:\the\dp0}
\setbox0=\vbox{\hsize=6.890cm\parindent=0pt\linespread{1}\fontsize{10}{13.50}\selectfont\bfseries\centering\noindent 范例\par}
\typeout{TCELL:p307_1:4:\the\ht0:\the\dp0}
\setbox0=\vbox{\hsize=2.440cm\parindent=0pt\linespread{1}\fontsize{10}{13.50}\selectfont\raggedright\noindent 一开始就不想付费\par}
\typeout{TCELL:p307_1:5:\the\ht0:\the\dp0}
\setbox0=\vbox{\hsize=1.540cm\parindent=0pt\linespread{1}\fontsize{10}{13.50}\selectfont\centering\noindent 有形财物\par}
\typeout{TCELL:p307_1:6:\the\ht0:\the\dp0}
\setbox0=\vbox{\hsize=1.540cm\parindent=0pt\linespread{1}\fontsize{10}{13.50}\selectfont\centering\noindent 欺骗手段\par}
\typeout{TCELL:p307_1:7:\the\ht0:\the\dp0}
\setbox0=\vbox{\hsize=2.290cm\parindent=0pt\linespread{1}\fontsize{10}{13.50}\selectfont\centering\noindent 构成诈骗罪\par}
\typeout{TCELL:p307_1:8:\the\ht0:\the\dp0}
\setbox0=\vbox{\hsize=6.890cm\parindent=0pt\linespread{1}\fontsize{10}{13.50}\selectfont\raggedright\noindent 甲一开始就不想付费，欺骗店家提供饭菜，吃完逃离。甲构成诈骗罪，行为对象是饭菜本身\par}
\typeout{TCELL:p307_1:9:\the\ht0:\the\dp0}
\setbox0=\vbox{\hsize=2.440cm\parindent=0pt\linespread{1}\fontsize{10}{13.50}\selectfont\raggedright\noindent 一开始想付费，产生债务后，不想付费\par}
\typeout{TCELL:p307_1:10:\the\ht0:\the\dp0}
\setbox0=\vbox{\hsize=1.540cm\parindent=0pt\linespread{1}\fontsize{10}{13.50}\selectfont\centering\noindent 财产性利益\par}
\typeout{TCELL:p307_1:11:\the\ht0:\the\dp0}
\setbox0=\vbox{\hsize=1.540cm\parindent=0pt\linespread{1}\fontsize{10}{13.50}\selectfont\centering\noindent 欺骗手段\par}
\typeout{TCELL:p307_1:12:\the\ht0:\the\dp0}
\setbox0=\vbox{\hsize=2.290cm\parindent=0pt\linespread{1}\fontsize{10}{13.50}\selectfont\centering\noindent 构成诈骗罪\par}
\typeout{TCELL:p307_1:13:\the\ht0:\the\dp0}
\setbox0=\vbox{\hsize=6.890cm\parindent=0pt\linespread{1}\fontsize{10}{13.50}\selectfont\raggedright\noindent 例1，甲开始有付费意思，吃完后用假币支付餐费。甲构成诈骗罪。\\{}例2，甲开始有付费意思，吃完后不想付费，欺骗老板：我下次来一起付。老板相信，放走甲。老板放弃了当场实现债权。甲构成诈骗罪\par}
\typeout{TCELL:p307_1:14:\the\ht0:\the\dp0}
\setbox0=\vbox{\hsize=1.540cm\parindent=0pt\linespread{1}\fontsize{10}{13.50}\selectfont\centering\noindent 平和手段\par}
\typeout{TCELL:p307_1:15:\the\ht0:\the\dp0}
\setbox0=\vbox{\hsize=2.290cm\parindent=0pt\linespread{1}\fontsize{10}{13.50}\selectfont\raggedright\noindent 多数说认为不构成盗窃罪。少数说认为构成盗窃罪\par}
\typeout{TCELL:p307_1:16:\the\ht0:\the\dp0}
\setbox0=\vbox{\hsize=6.890cm\parindent=0pt\linespread{1}\fontsize{10}{13.50}\selectfont\raggedright\noindent 例1，甲开始有付费意思，吃完后不想付钱，从后门偷偷溜走。存在观点展示。\\{}例2，甲开始有付费意思，吃完后不想付钱，欺骗老板：我送下朋友就回来结账，然后溜走。老板没有放弃当场实现债权。存在观点展示\par}
\typeout{TCELL:p307_1:17:\the\ht0:\the\dp0}
\setbox0=\vbox{\hsize=2.440cm\parindent=0pt\linespread{1}\fontsize{10}{13.50}\selectfont\bfseries\centering\noindent 总结“掐卡案”\par}
\typeout{TCELL:p311_1:0:\the\ht0:\the\dp0}
\setbox0=\vbox{\hsize=3.940cm\parindent=0pt\linespread{1}\fontsize{10}{13.50}\selectfont\bfseries\centering\noindent 观点展示\par}
\typeout{TCELL:p311_1:1:\the\ht0:\the\dp0}
\setbox0=\vbox{\hsize=1.290cm\parindent=0pt\linespread{1}\fontsize{10}{13.50}\selectfont\bfseries\centering\noindent 掐卡\par}
\typeout{TCELL:p311_1:2:\the\ht0:\the\dp0}
\setbox0=\vbox{\hsize=2.590cm\parindent=0pt\linespread{1}\fontsize{10}{13.50}\selectfont\bfseries\centering\noindent 转账\par}
\typeout{TCELL:p311_1:3:\the\ht0:\the\dp0}
\setbox0=\vbox{\hsize=4.440cm\parindent=0pt\linespread{1}\fontsize{10}{13.50}\selectfont\bfseries\centering\noindent 取款（取银行现金）\par}
\typeout{TCELL:p311_1:4:\the\ht0:\the\dp0}
\setbox0=\vbox{\hsize=2.440cm\parindent=0pt\linespread{1}\fontsize{10}{13.50}\selectfont\centering\noindent 卖卡后，针对合法资金\par}
\typeout{TCELL:p311_1:5:\the\ht0:\the\dp0}
\setbox0=\vbox{\hsize=3.940cm\parindent=0pt\linespread{1}\fontsize{10}{13.50}\selectfont\centering\noindent 事实占有说（多数说）\par}
\typeout{TCELL:p311_1:6:\the\ht0:\the\dp0}
\setbox0=\vbox{\hsize=1.290cm\parindent=0pt\linespread{1}\fontsize{10}{13.50}\selectfont\centering\noindent 盗窃罪\par}
\typeout{TCELL:p311_1:7:\the\ht0:\the\dp0}
\setbox0=\vbox{\hsize=2.590cm\parindent=0pt\linespread{1}\fontsize{10}{13.50}\selectfont\raggedright\noindent 掐卡定了盗窃，转账不再定罪\par}
\typeout{TCELL:p311_1:8:\the\ht0:\the\dp0}
\setbox0=\vbox{\hsize=4.440cm\parindent=0pt\linespread{1}\fontsize{10}{13.50}\selectfont\raggedright\noindent 多数说：对银行构成犯罪。\\{}少数说：对银行不构成犯罪\par}
\typeout{TCELL:p311_1:9:\the\ht0:\the\dp0}
\setbox0=\vbox{\hsize=3.940cm\parindent=0pt\linespread{1}\fontsize{10}{13.50}\selectfont\centering\noindent 法律占有说（少数说）\par}
\typeout{TCELL:p311_1:10:\the\ht0:\the\dp0}
\setbox0=\vbox{\hsize=8.840cm\parindent=0pt\linespread{1}\fontsize{10}{13.50}\selectfont\centering\noindent 侵占罪\par}
\typeout{TCELL:p311_1:11:\the\ht0:\the\dp0}
\setbox0=\vbox{\hsize=2.440cm\parindent=0pt\linespread{1}\fontsize{10}{13.50}\selectfont\centering\noindent 卖卡后，针对赃款\par}
\typeout{TCELL:p311_1:12:\the\ht0:\the\dp0}
\setbox0=\vbox{\hsize=3.940cm\parindent=0pt\linespread{1}\fontsize{10}{13.50}\selectfont\centering\noindent 事实占有说（多数说）\par}
\typeout{TCELL:p311_1:13:\the\ht0:\the\dp0}
\setbox0=\vbox{\hsize=8.840cm\parindent=0pt\linespread{1}\fontsize{10}{13.50}\selectfont\centering\noindent 盗窃罪\par}
\typeout{TCELL:p311_1:14:\the\ht0:\the\dp0}
\setbox0=\vbox{\hsize=3.940cm\parindent=0pt\linespread{1}\fontsize{10}{13.50}\selectfont\centering\noindent 法律占有说（少数说）\par}
\typeout{TCELL:p311_1:15:\the\ht0:\the\dp0}
\setbox0=\vbox{\hsize=8.840cm\parindent=0pt\linespread{1}\fontsize{10}{13.50}\selectfont\centering\noindent 侵占罪\par}
\typeout{TCELL:p311_1:16:\the\ht0:\the\dp0}
\setbox0=\vbox{\hsize=3.840cm\parindent=0pt\linespread{1}\fontsize{10}{13.50}\selectfont\bfseries\centering\noindent 四罪与故意杀人罪的关系\par}
\typeout{TCELL:p318_1:0:\the\ht0:\the\dp0}
\setbox0=\vbox{\hsize=3.040cm\parindent=0pt\linespread{1}\fontsize{10}{13.50}\selectfont\bfseries\centering\noindent 结论\par}
\typeout{TCELL:p318_1:1:\the\ht0:\the\dp0}
\setbox0=\vbox{\hsize=8.340cm\parindent=0pt\linespread{1}\fontsize{10}{13.50}\selectfont\bfseries\centering\noindent 范例\par}
\typeout{TCELL:p318_1:2:\the\ht0:\the\dp0}
\setbox0=\vbox{\hsize=3.840cm\parindent=0pt\linespread{1}\fontsize{10}{13.50}\selectfont\centering\noindent 实施放火，危害了公共安全，对死亡结果持过失心理\par}
\typeout{TCELL:p318_1:3:\the\ht0:\the\dp0}
\setbox0=\vbox{\hsize=3.040cm\parindent=0pt\linespread{1}\fontsize{10}{13.50}\selectfont\raggedright\noindent 成立结果加重犯：放火罪（过失）致人死亡\par}
\typeout{TCELL:p318_1:4:\the\ht0:\the\dp0}
\setbox0=\vbox{\hsize=8.340cm\parindent=0pt\linespread{1}\fontsize{10}{13.50}\selectfont\raggedright\noindent 甲欲报复单位，认为单位仓库没人把守，便放火烧仓库，未料火灾将睡着的库管人员烧死。甲构成放火罪（过失）致人死亡\par}
\typeout{TCELL:p318_1:5:\the\ht0:\the\dp0}
\setbox0=\vbox{\hsize=3.840cm\parindent=0pt\linespread{1}\fontsize{10}{13.50}\selectfont\centering\noindent 实施放火，危害了公共安全，对死亡结果持故意心理\par}
\typeout{TCELL:p318_1:6:\the\ht0:\the\dp0}
\setbox0=\vbox{\hsize=3.040cm\parindent=0pt\linespread{1}\fontsize{10}{13.50}\selectfont\raggedright\noindent 成立结果加重犯：放火罪（故意）致人死亡\par}
\typeout{TCELL:p318_1:7:\the\ht0:\the\dp0}
\setbox0=\vbox{\hsize=8.340cm\parindent=0pt\linespread{1}\fontsize{10}{13.50}\selectfont\raggedright\noindent 甲欲报复社会，在公交车上放火，烧死多人。甲构成放火罪（故意）致人死亡①\par}
\typeout{TCELL:p318_1:8:\the\ht0:\the\dp0}
\setbox0=\vbox{\hsize=3.840cm\parindent=0pt\linespread{1}\fontsize{10}{13.50}\selectfont\centering\noindent 实施放火，没有危害公共安全，对死亡结果持故意心理\par}
\typeout{TCELL:p318_1:9:\the\ht0:\the\dp0}
\setbox0=\vbox{\hsize=3.040cm\parindent=0pt\linespread{1}\fontsize{10}{13.50}\selectfont\raggedright\noindent 只构成故意杀人罪\par}
\typeout{TCELL:p318_1:10:\the\ht0:\the\dp0}
\setbox0=\vbox{\hsize=8.340cm\parindent=0pt\linespread{1}\fontsize{10}{13.50}\selectfont\raggedright\noindent 甲欲杀死乙全家，放火烧乙家住宅，乙一家三口全部死亡。乙家住宅是独门独户，四周无人。甲的行为没有危害公共安全，只构成故意杀人罪\par}
\typeout{TCELL:p318_1:11:\the\ht0:\the\dp0}
\setbox0=\vbox{\hsize=3.590cm\parindent=0pt\linespread{1}\fontsize{10}{13.50}\selectfont\bfseries\centering\noindent 以危险方法危害公共安全罪的构成要件\par}
\typeout{TCELL:p319_1:0:\the\ht0:\the\dp0}
\setbox0=\vbox{\hsize=5.790cm\parindent=0pt\linespread{1}\fontsize{10}{13.50}\selectfont\bfseries\centering\noindent 要点\par}
\typeout{TCELL:p319_1:1:\the\ht0:\the\dp0}
\setbox0=\vbox{\hsize=5.840cm\parindent=0pt\linespread{1}\fontsize{10}{13.50}\selectfont\bfseries\centering\noindent 范例\par}
\typeout{TCELL:p319_1:2:\the\ht0:\the\dp0}
\setbox0=\vbox{\hsize=3.590cm\parindent=0pt\linespread{1}\fontsize{10}{13.50}\selectfont\centering\noindent 危险的性质\par}
\typeout{TCELL:p319_1:3:\the\ht0:\the\dp0}
\setbox0=\vbox{\hsize=5.790cm\parindent=0pt\linespread{1}\fontsize{10}{13.50}\selectfont\raggedright\noindent 必须是具体危险，不能是抽象危险\par}
\typeout{TCELL:p319_1:4:\the\ht0:\the\dp0}
\setbox0=\vbox{\hsize=5.840cm\parindent=0pt\linespread{1}\fontsize{10}{13.50}\selectfont\raggedright\noindent 甲盗窃公共草坪上的消防栓铜芯。法院判决甲构成以危险方法危害公共安全罪。该判决有误。盗窃行为对公共安全只产生抽象危险。甲只构成盗窃罪（2023年试题）\par}
\typeout{TCELL:p319_1:5:\the\ht0:\the\dp0}
\setbox0=\vbox{\hsize=3.590cm\parindent=0pt\linespread{1}\fontsize{10}{13.50}\selectfont\centering\noindent 危险的内容\par}
\typeout{TCELL:p319_1:6:\the\ht0:\the\dp0}
\setbox0=\vbox{\hsize=5.790cm\parindent=0pt\linespread{1}\fontsize{10}{13.50}\selectfont\raggedright\noindent 限于人身伤亡和财产损失，不包括精神损害，也即限于物质性损害，不包括非物质性损害\par}
\typeout{TCELL:p319_1:7:\the\ht0:\the\dp0}
\setbox0=\vbox{\hsize=5.840cm\parindent=0pt\linespread{1}\fontsize{10}{13.50}\selectfont\raggedright\noindent 甲将面粉装入信封内，写上“炭疽热病菌粉”，寄给乙单位，引起乙单位恐慌。法院判决甲构成以危险方法危害公共安全罪。该判决有误。甲没有危害公共安全。甲应构成投放虚假危险物质罪，该罪的保护法益不是公共安全，而是社会管理秩序\par}
\typeout{TCELL:p319_1:8:\the\ht0:\the\dp0}
\setbox0=\vbox{\hsize=3.590cm\parindent=0pt\linespread{1}\fontsize{10}{13.50}\selectfont\centering\noindent 危险的样态\par}
\typeout{TCELL:p319_1:9:\the\ht0:\the\dp0}
\setbox0=\vbox{\hsize=5.790cm\parindent=0pt\linespread{1}\fontsize{10}{13.50}\selectfont\raggedright\noindent 公共危险是指危险\CJKunderwave{具有危及多数人的}\CJKunderwave{可能性}\par}
\typeout{TCELL:p319_1:10:\the\ht0:\the\dp0}
\setbox0=\vbox{\hsize=5.840cm\parindent=0pt\linespread{1}\fontsize{10}{13.50}\selectfont\raggedright\noindent 甲用刀在小学门口刺伤许多学生。法院判决甲构成以危险方法危害公共安全罪。该判决有误。甲应构成数个故意伤害罪\par}
\typeout{TCELL:p319_1:11:\the\ht0:\the\dp0}
\setbox0=\vbox{\hsize=3.590cm\parindent=0pt\linespread{1}\fontsize{10}{13.50}\selectfont\centering\noindent 危险的程度\par}
\typeout{TCELL:p319_1:12:\the\ht0:\the\dp0}
\setbox0=\vbox{\hsize=5.790cm\parindent=0pt\linespread{1}\fontsize{10}{13.50}\selectfont\raggedright\noindent 危险的程度要达到与放火罪、爆炸罪、投放危险物质罪的相同程度。\\{}本罪被称为兜底罪名。错误思维是，如果行为达不到放火罪的程度，那就降一格，适用兜底罪名。实际上，本罪与放火罪是并列关系，而非高低关系\par}
\typeout{TCELL:p319_1:13:\the\ht0:\the\dp0}
\setbox0=\vbox{\hsize=5.840cm\parindent=0pt\linespread{1}\fontsize{10}{13.50}\selectfont\raggedright\noindent 甲为了报复社会，在高速公路上点燃一个很小的火堆。法院判决，危险程度不高，不定放火罪，但定以危险方法危害公共安全罪。该判决有误。危险既然没达到放火罪的程度，也就没达到本罪程度，不构成本罪\par}
\typeout{TCELL:p319_1:14:\the\ht0:\the\dp0}
\setbox0=\vbox{\hsize=3.590cm\parindent=0pt\linespread{1}\fontsize{10}{13.50}\selectfont\centering\noindent 适用顺序\par}
\typeout{TCELL:p319_1:15:\the\ht0:\the\dp0}
\setbox0=\vbox{\hsize=5.790cm\parindent=0pt\linespread{1}\fontsize{10}{13.50}\selectfont\raggedright\noindent 本罪实际上是补充罪名。如果行为符合其他犯罪的构成要件，应优先定其他犯罪\par}
\typeout{TCELL:p319_1:16:\the\ht0:\the\dp0}
\setbox0=\vbox{\hsize=5.840cm\parindent=0pt\linespread{1}\fontsize{10}{13.50}\selectfont\raggedright\noindent 例1，劫持火车、电车的行为，应定破坏交通工具罪。\\{}例2，盗窃机动车道路上的井盖，危害公共安全，应定破坏交通设施罪（2020年、2025年试题）\par}
\typeout{TCELL:p319_1:17:\the\ht0:\the\dp0}
\setbox0=\vbox{\hsize=1.740cm\parindent=0pt\linespread{1}\fontsize{10}{13.50}\selectfont\bfseries\centering\noindent 构成要件\par}
\typeout{TCELL:p322_1:0:\the\ht0:\the\dp0}
\setbox0=\vbox{\hsize=3.790cm\parindent=0pt\linespread{1}\fontsize{10}{13.50}\selectfont\bfseries\centering\noindent 要点\par}
\typeout{TCELL:p322_1:1:\the\ht0:\the\dp0}
\setbox0=\vbox{\hsize=9.690cm\parindent=0pt\linespread{1}\fontsize{10}{13.50}\selectfont\bfseries\centering\noindent 范例\par}
\typeout{TCELL:p322_1:2:\the\ht0:\the\dp0}
\setbox0=\vbox{\hsize=1.740cm\parindent=0pt\linespread{1}\fontsize{10}{13.50}\selectfont\centering\noindent 行为主体\par}
\typeout{TCELL:p322_1:3:\the\ht0:\the\dp0}
\setbox0=\vbox{\hsize=3.790cm\parindent=0pt\linespread{1}\fontsize{10}{13.50}\selectfont\raggedright\noindent 不限于驾驶者\par}
\typeout{TCELL:p322_1:4:\the\ht0:\the\dp0}
\setbox0=\vbox{\hsize=9.690cm\parindent=0pt\linespread{1}\fontsize{10}{13.50}\selectfont\raggedright\noindent 甲看到乙醉酒，仍将车交给乙驾驶，乙醉驾不慎撞死人。乙构成交通肇事罪，甲也构成交通肇事罪\par}
\typeout{TCELL:p322_1:5:\the\ht0:\the\dp0}
\setbox0=\vbox{\hsize=1.740cm\parindent=0pt\linespread{1}\fontsize{10}{13.50}\selectfont\centering\noindent 行为地点\par}
\typeout{TCELL:p322_1:6:\the\ht0:\the\dp0}
\setbox0=\vbox{\hsize=3.790cm\parindent=0pt\linespread{1}\fontsize{10}{13.50}\selectfont\raggedright\noindent 公共交通领域\par}
\typeout{TCELL:p322_1:7:\the\ht0:\the\dp0}
\setbox0=\vbox{\hsize=9.690cm\parindent=0pt\linespread{1}\fontsize{10}{13.50}\selectfont\raggedright\noindent （1）在公共交通领域外，驾驶车辆过失致人死亡，构成过失致人死亡罪。\\{}(2）公共交通领域包括公路交通、水上交通。（2025年试题）\\{}（3）铁路事故，定铁路运营安全事故罪。\\{}（4）飞行事故，定重大飞行事故罪\par}
\typeout{TCELL:p322_1:8:\the\ht0:\the\dp0}
\setbox0=\vbox{\hsize=1.740cm\parindent=0pt\linespread{1}\fontsize{10}{13.50}\selectfont\centering\noindent 肇事行为\par}
\typeout{TCELL:p322_1:9:\the\ht0:\the\dp0}
\setbox0=\vbox{\hsize=3.790cm\parindent=0pt\linespread{1}\fontsize{10}{13.50}\selectfont\raggedright\noindent 肇事行为是过失行为\par}
\typeout{TCELL:p322_1:10:\the\ht0:\the\dp0}
\setbox0=\vbox{\hsize=9.690cm\parindent=0pt\linespread{1}\fontsize{10}{13.50}\selectfont\raggedright\noindent 甲驾车想闯红灯，认为不会出事，不慎撞死行人。甲闯红灯是故意为之（行政法上的故意），但对死亡结果是过失心态（刑法上的过失)\par}
\typeout{TCELL:p322_1:11:\the\ht0:\the\dp0}
\setbox0=\vbox{\hsize=1.740cm\parindent=0pt\linespread{1}\fontsize{10}{13.50}\selectfont\centering\noindent 实害结果\par}
\typeout{TCELL:p322_1:12:\the\ht0:\the\dp0}
\setbox0=\vbox{\hsize=3.790cm\parindent=0pt\linespread{1}\fontsize{10}{13.50}\selectfont\raggedright\noindent 死亡1人；重伤3人；重伤1人，并有严重情节（酒驾、吸毒驾驶、无照驾驶、严重超载、肇事后逃逸)\par}
\typeout{TCELL:p322_1:13:\the\ht0:\the\dp0}
\setbox0=\vbox{\hsize=9.690cm\parindent=0pt\linespread{1}\fontsize{10}{13.50}\selectfont\raggedright\noindent 这里的“肇事后逃逸”被作为本罪成立条件后，便不能再作为本罪的法定刑升格条件“肇事后逃逸，处3年至7年有期徒刑”，否则属于重复评价、重复处罚\par}
\typeout{TCELL:p322_1:14:\the\ht0:\the\dp0}
\setbox0=\vbox{\hsize=1.740cm\parindent=0pt\linespread{1}\fontsize{10}{13.50}\selectfont\centering\noindent 因果关系\par}
\typeout{TCELL:p322_1:15:\the\ht0:\the\dp0}
\setbox0=\vbox{\hsize=3.790cm\parindent=0pt\linespread{1}\fontsize{10}{13.50}\selectfont\raggedright\noindent 运用因果关系原理。特别是，实害结果必须是规范保护范围内的结果\par}
\typeout{TCELL:p322_1:16:\the\ht0:\the\dp0}
\setbox0=\vbox{\hsize=9.690cm\parindent=0pt\linespread{1}\fontsize{10}{13.50}\selectfont\raggedright\noindent 甲驾驶没有年检的车辆，但该车并无故障。乙横穿高速公路，甲刹车不及撞死乙。“禁止驾驶未年检的车辆”，这个规范要防止的结果是，因车辆故障导致交通事故。现在乙死亡这个结果不是这个规范所能防止的结果，不能归属给甲。甲不构成交通肇事罪。\\{}[注意］不能将行政法上的违法直接等同于刑事违法，不能将行政法上的责任直接等同于刑事责任\par}
\typeout{TCELL:p322_1:17:\the\ht0:\the\dp0}
\setbox0=\vbox{\hsize=3.590cm\parindent=0pt\linespread{1}\fontsize{10}{13.50}\selectfont\bfseries\centering\noindent 罪数关系\par}
\typeout{TCELL:p327_1:0:\the\ht0:\the\dp0}
\setbox0=\vbox{\hsize=3.590cm\parindent=0pt\linespread{1}\fontsize{10}{13.50}\selectfont\bfseries\centering\noindent 行为类型\par}
\typeout{TCELL:p327_1:1:\the\ht0:\the\dp0}
\setbox0=\vbox{\hsize=8.040cm\parindent=0pt\linespread{1}\fontsize{10}{13.50}\selectfont\bfseries\centering\noindent 范例\par}
\typeout{TCELL:p327_1:2:\the\ht0:\the\dp0}
\setbox0=\vbox{\hsize=3.590cm\parindent=0pt\linespread{1}\fontsize{10}{13.50}\selectfont\centering\noindent 危险驾驶罪与交通肇事罪\par}
\typeout{TCELL:p327_1:3:\the\ht0:\the\dp0}
\setbox0=\vbox{\hsize=3.590cm\parindent=0pt\linespread{1}\fontsize{10}{13.50}\selectfont\centering\noindent 一个危险流\par}
\typeout{TCELL:p327_1:4:\the\ht0:\the\dp0}
\setbox0=\vbox{\hsize=8.040cm\parindent=0pt\linespread{1}\fontsize{10}{13.50}\selectfont\raggedright\noindent 甲醉驾行驶了10公里，在第11公里因醉驾而过失撞死行人。第一阶段，甲已构成危险驾驶罪。第二阶段，甲又构成交通肇事罪。由于是同一个危险流（醉驾制造的危险流），重罪吸收轻罪，最终只定交通肇事罪\par}
\typeout{TCELL:p327_1:5:\the\ht0:\the\dp0}
\setbox0=\vbox{\hsize=3.590cm\parindent=0pt\linespread{1}\fontsize{10}{13.50}\selectfont\centering\noindent 两个危险流\par}
\typeout{TCELL:p327_1:6:\the\ht0:\the\dp0}
\setbox0=\vbox{\hsize=8.040cm\parindent=0pt\linespread{1}\fontsize{10}{13.50}\selectfont\raggedright\noindent 甲醉驾行驶了10公里，在第11公里，不是因为醉驾，而是因为新的违章行为（如闯红灯）过失撞死行人。第一阶段，甲已构成危险驾驶罪。第二阶段，甲又构成交通肇事罪。由于是两个不同的危险流（醉驾制造的危险流和闯红灯制造的危险流），应数罪并罚\par}
\typeout{TCELL:p327_1:7:\the\ht0:\the\dp0}
\setbox0=\vbox{\hsize=3.590cm\parindent=0pt\linespread{1}\fontsize{10}{13.50}\selectfont\centering\noindent 危险驾驶罪与以危险方法危害公共安全罪\par}
\typeout{TCELL:p327_1:8:\the\ht0:\the\dp0}
\setbox0=\vbox{\hsize=3.590cm\parindent=0pt\linespread{1}\fontsize{10}{13.50}\selectfont\centering\noindent 一个行为\par}
\typeout{TCELL:p327_1:9:\the\ht0:\the\dp0}
\setbox0=\vbox{\hsize=8.040cm\parindent=0pt\linespread{1}\fontsize{10}{13.50}\selectfont\raggedright\noindent 甲醉酒在高速公路上逆向行驶。甲的一个行为同时触犯危险驾驶罪和以危险方法危害公共安全罪，想象竞合，择一重罪论处，定以危险方法危害公共安全罪\par}
\typeout{TCELL:p327_1:10:\the\ht0:\the\dp0}
\setbox0=\vbox{\hsize=3.590cm\parindent=0pt\linespread{1}\fontsize{10}{13.50}\selectfont\centering\noindent 两个行为\par}
\typeout{TCELL:p327_1:11:\the\ht0:\the\dp0}
\setbox0=\vbox{\hsize=8.040cm\parindent=0pt\linespread{1}\fontsize{10}{13.50}\selectfont\raggedright\noindent 甲醉驾行驶了5公里，被警察发现，为了逃跑，横冲直撞，撞死行人。甲的第一个行为构成危险驾驶罪，第二个行为构成以危险方法危害公共安全罪，数罪并罚\par}
\typeout{TCELL:p327_1:12:\the\ht0:\the\dp0}
\setbox0=\vbox{\hsize=3.590cm\parindent=0pt\linespread{1}\fontsize{10}{13.50}\selectfont\centering\noindent 交通肇事罪与以危险方法危害公共安全罪\par}
\typeout{TCELL:p327_1:13:\the\ht0:\the\dp0}
\setbox0=\vbox{\hsize=3.590cm\parindent=0pt\linespread{1}\fontsize{10}{13.50}\selectfont\raggedright\noindent 先构成交通肇事罪，后构成以危险方法危害公共安全罪\par}
\typeout{TCELL:p327_1:14:\the\ht0:\the\dp0}
\setbox0=\vbox{\hsize=8.040cm\parindent=0pt\linespread{1}\fontsize{10}{13.50}\selectfont\raggedright\noindent 并罚。例如，甲先构成交通肇事罪，后逃逸中横冲直撞，构成以危险方法危害公共安全罪。两罪并罚\par}
\typeout{TCELL:p327_1:15:\the\ht0:\the\dp0}
\setbox0=\vbox{\hsize=3.590cm\parindent=0pt\linespread{1}\fontsize{10}{13.50}\selectfont\raggedright\noindent 先构成以危险方法危害公共安全罪，后构成交通肇事罪\par}
\typeout{TCELL:p327_1:16:\the\ht0:\the\dp0}
\setbox0=\vbox{\hsize=8.040cm\parindent=0pt\linespread{1}\fontsize{10}{13.50}\selectfont\raggedright\noindent 吸收。例如，甲驾车在高速公路上高速逆行，10公里后不慎撞上汽车，导致对方死亡。这是一个危险流演变为实害结果。以危险方法危害公共安全罪吸收交通肇事罪\par}
\typeout{TCELL:p327_1:17:\the\ht0:\the\dp0}
\setbox0=\vbox{\hsize=3.590cm\parindent=0pt\linespread{1}\fontsize{10}{13.50}\selectfont\centering\noindent 危险驾驶罪与运输类犯罪（运输假币罪，运输毒品罪）\par}
\typeout{TCELL:p327_1:18:\the\ht0:\the\dp0}
\setbox0=\vbox{\hsize=3.590cm\parindent=0pt\linespread{1}\fontsize{10}{13.50}\selectfont\centering\noindent 一个行为\par}
\typeout{TCELL:p327_1:19:\the\ht0:\the\dp0}
\setbox0=\vbox{\hsize=8.040cm\parindent=0pt\linespread{1}\fontsize{10}{13.50}\selectfont\raggedright\noindent 甲醉酒驾车运输毒品，一个驾驶运输行为同时触犯危险驾驶罪和运输毒品罪，想象竞合，择一重罪论处，定运输毒品罪\par}
\typeout{TCELL:p327_1:20:\the\ht0:\the\dp0}
\setbox0=\vbox{\hsize=1.740cm\parindent=0pt\linespread{1}\fontsize{10}{13.50}\selectfont\bfseries\centering\noindent 主观心理\par}
\typeout{TCELL:p328_1:0:\the\ht0:\the\dp0}
\setbox0=\vbox{\hsize=6.940cm\parindent=0pt\linespread{1}\fontsize{10}{13.50}\selectfont\bfseries\centering\noindent 罪名\par}
\typeout{TCELL:p328_1:1:\the\ht0:\the\dp0}
\setbox0=\vbox{\hsize=6.540cm\parindent=0pt\linespread{1}\fontsize{10}{13.50}\selectfont\bfseries\centering\noindent 范例\par}
\typeout{TCELL:p328_1:2:\the\ht0:\the\dp0}
\setbox0=\vbox{\hsize=1.740cm\parindent=0pt\linespread{1}\fontsize{10}{13.50}\selectfont\centering\noindent 故意犯罪\par}
\typeout{TCELL:p328_1:3:\the\ht0:\the\dp0}
\setbox0=\vbox{\hsize=6.940cm\parindent=0pt\linespread{1}\fontsize{10}{13.50}\selectfont\raggedright\noindent 危险达到放火罪、爆炸罪的相当程度：以危险方法危害公共安全罪\par}
\typeout{TCELL:p328_1:4:\the\ht0:\the\dp0}
\setbox0=\vbox{\hsize=6.540cm\parindent=0pt\linespread{1}\fontsize{10}{13.50}\selectfont\raggedright\noindent 甲在高速公路上故意逆向高速行驶。甲构成以危险方法危害公共安全罪\par}
\typeout{TCELL:p328_1:5:\the\ht0:\the\dp0}
\setbox0=\vbox{\hsize=6.940cm\parindent=0pt\linespread{1}\fontsize{10}{13.50}\selectfont\raggedright\noindent 危险未达到放火罪、爆炸罪的相当程度：危险驾驶罪、妨害安全驾驶罪\par}
\typeout{TCELL:p328_1:6:\the\ht0:\the\dp0}
\setbox0=\vbox{\hsize=6.540cm\parindent=0pt\linespread{1}\fontsize{10}{13.50}\selectfont\raggedright\noindent 例1，甲在道路上醉酒驾驶，构成危险驾驶罪。\\{}例2，乘客短时间内抢夺公交车的方向盘，构成妨害安全驾驶罪\par}
\typeout{TCELL:p328_1:7:\the\ht0:\the\dp0}
\setbox0=\vbox{\hsize=1.740cm\parindent=0pt\linespread{1}\fontsize{10}{13.50}\selectfont\centering\noindent 过失犯罪\par}
\typeout{TCELL:p328_1:8:\the\ht0:\the\dp0}
\setbox0=\vbox{\hsize=6.940cm\parindent=0pt\linespread{1}\fontsize{10}{13.50}\selectfont\raggedright\noindent 具体罪名：交通肇事罪。\\{}兜底罪名（补充罪名）：过失以危险方法危害公共安全罪\par}
\typeout{TCELL:p328_1:9:\the\ht0:\the\dp0}
\setbox0=\vbox{\hsize=6.540cm\parindent=0pt\linespread{1}\fontsize{10}{13.50}\selectfont\raggedright\noindent 甲驾车闯红灯，不慎撞死人，构成交通肇事罪\par}
\typeout{TCELL:p328_1:10:\the\ht0:\the\dp0}
\setbox0=\vbox{\hsize=2.890cm\parindent=0pt\linespread{1}\fontsize{10}{13.50}\selectfont\bfseries\centering\noindent 领域\par}
\typeout{TCELL:p332_1:0:\the\ht0:\the\dp0}
\setbox0=\vbox{\hsize=5.290cm\parindent=0pt\linespread{1}\fontsize{10}{13.50}\selectfont\bfseries\centering\noindent 危险\par}
\typeout{TCELL:p332_1:1:\the\ht0:\the\dp0}
\setbox0=\vbox{\hsize=7.040cm\parindent=0pt\linespread{1}\fontsize{10}{13.50}\selectfont\bfseries\centering\noindent 实害结果\par}
\typeout{TCELL:p332_1:2:\the\ht0:\the\dp0}
\setbox0=\vbox{\hsize=2.890cm\parindent=0pt\linespread{1}\fontsize{10}{13.50}\selectfont\centering\noindent 生产领域\par}
\typeout{TCELL:p332_1:3:\the\ht0:\the\dp0}
\setbox0=\vbox{\hsize=5.290cm\parindent=0pt\linespread{1}\fontsize{10}{13.50}\selectfont\raggedright\noindent 危险作业罪（故意犯罪）\par}
\typeout{TCELL:p332_1:4:\the\ht0:\the\dp0}
\setbox0=\vbox{\hsize=7.040cm\parindent=0pt\linespread{1}\fontsize{10}{13.50}\selectfont\raggedright\noindent 重大责任事故罪（过失犯罪）（2025年试题）\par}
\typeout{TCELL:p332_1:5:\the\ht0:\the\dp0}
\setbox0=\vbox{\hsize=2.890cm\parindent=0pt\linespread{1}\fontsize{10}{13.50}\selectfont\centering\noindent 交通领域\par}
\typeout{TCELL:p332_1:6:\the\ht0:\the\dp0}
\setbox0=\vbox{\hsize=5.290cm\parindent=0pt\linespread{1}\fontsize{10}{13.50}\selectfont\raggedright\noindent 危险驾驶罪（故意犯罪）\par}
\typeout{TCELL:p332_1:7:\the\ht0:\the\dp0}
\setbox0=\vbox{\hsize=7.040cm\parindent=0pt\linespread{1}\fontsize{10}{13.50}\selectfont\raggedright\noindent 交通肇事罪（过失犯罪）\par}
\typeout{TCELL:p332_1:8:\the\ht0:\the\dp0}
\setbox0=\vbox{\hsize=2.190cm\parindent=0pt\linespread{1}\fontsize{10}{13.50}\selectfont\bfseries\centering\noindent 食品犯罪\par}
\typeout{TCELL:p334_1:0:\the\ht0:\the\dp0}
\setbox0=\vbox{\hsize=2.140cm\parindent=0pt\linespread{1}\fontsize{10}{13.50}\selectfont\bfseries\centering\noindent 成立条件\par}
\typeout{TCELL:p334_1:1:\the\ht0:\the\dp0}
\setbox0=\vbox{\hsize=2.390cm\parindent=0pt\linespread{1}\fontsize{10}{13.50}\selectfont\bfseries\centering\noindent 食品类型\par}
\typeout{TCELL:p334_1:2:\the\ht0:\the\dp0}
\setbox0=\vbox{\hsize=8.240cm\parindent=0pt\linespread{1}\fontsize{10}{13.50}\selectfont\bfseries\centering\noindent 范例\par}
\typeout{TCELL:p334_1:3:\the\ht0:\the\dp0}
\setbox0=\vbox{\hsize=2.190cm\parindent=0pt\linespread{1}\fontsize{10}{13.50}\selectfont\centering\noindent 生产、销售有毒、有害食品罪\par}
\typeout{TCELL:p334_1:4:\the\ht0:\the\dp0}
\setbox0=\vbox{\hsize=2.140cm\parindent=0pt\linespread{1}\fontsize{10}{13.50}\selectfont\centering\noindent 抽象危险犯（也即行为犯）\par}
\typeout{TCELL:p334_1:5:\the\ht0:\the\dp0}
\setbox0=\vbox{\hsize=2.390cm\parindent=0pt\linespread{1}\fontsize{10}{13.50}\selectfont\raggedright\noindent 有毒、有害食品，是指掺入有毒、有害的非食品原料的食品\par}
\typeout{TCELL:p334_1:6:\the\ht0:\the\dp0}
\setbox0=\vbox{\hsize=8.240cm\parindent=0pt\linespread{1}\fontsize{10}{13.50}\selectfont\raggedright\noindent 1.用“瘦肉精”喂猪或销售这样的猪肉。注意：“食品”包括食材，如生猪、活鱼、生长中的蘑菇等。（2022年试题）\\{}2.将工业酒精勾兑成散装白酒出售。（2022年试题）\\{}3.将工业用猪油冒充食用油出售。\\{}4.用“地沟油”生产食用油\par}
\typeout{TCELL:p334_1:7:\the\ht0:\the\dp0}
\setbox0=\vbox{\hsize=2.190cm\parindent=0pt\linespread{1}\fontsize{10}{13.50}\selectfont\centering\noindent 生产、销售不符合安全标准的食品罪\par}
\typeout{TCELL:p334_1:8:\the\ht0:\the\dp0}
\setbox0=\vbox{\hsize=2.140cm\parindent=0pt\linespread{1}\fontsize{10}{13.50}\selectfont\centering\noindent 具体危险犯\par}
\typeout{TCELL:p334_1:9:\the\ht0:\the\dp0}
\setbox0=\vbox{\hsize=2.390cm\parindent=0pt\linespread{1}\fontsize{10}{13.50}\selectfont\raggedright\noindent 不符合安全标准的食品，没有达到毒害性\par}
\typeout{TCELL:p334_1:10:\the\ht0:\the\dp0}
\setbox0=\vbox{\hsize=8.240cm\parindent=0pt\linespread{1}\fontsize{10}{13.50}\selectfont\raggedright\noindent 1.含有严重超出标准限量的农药残留、重金属、污染物质的食品。\\{}2.包装材料、容器被污染的食品。\\{}3.超过保质期的食品\par}
\typeout{TCELL:p334_1:11:\the\ht0:\the\dp0}
\setbox0=\vbox{\hsize=1.990cm\parindent=0pt\linespread{1}\fontsize{10}{13.50}\selectfont\bfseries\centering\noindent 罪名\par}
\typeout{TCELL:p335_1:0:\the\ht0:\the\dp0}
\setbox0=\vbox{\hsize=1.990cm\parindent=0pt\linespread{1}\fontsize{10}{13.50}\selectfont\bfseries\centering\noindent 成立条件\par}
\typeout{TCELL:p335_1:1:\the\ht0:\the\dp0}
\setbox0=\vbox{\hsize=6.890cm\parindent=0pt\linespread{1}\fontsize{10}{13.50}\selectfont\bfseries\centering\noindent 行为对象和方式\par}
\typeout{TCELL:p335_1:2:\the\ht0:\the\dp0}
\setbox0=\vbox{\hsize=4.090cm\parindent=0pt\linespread{1}\fontsize{10}{13.50}\selectfont\bfseries\centering\noindent 要点\par}
\typeout{TCELL:p335_1:3:\the\ht0:\the\dp0}
\setbox0=\vbox{\hsize=1.990cm\parindent=0pt\linespread{1}\fontsize{10}{13.50}\selectfont\centering\noindent 生产、销售、提供假药罪\par}
\typeout{TCELL:p335_1:4:\the\ht0:\the\dp0}
\setbox0=\vbox{\hsize=1.990cm\parindent=0pt\linespread{1}\fontsize{10}{13.50}\selectfont\centering\noindent 抽象危险犯（行为犯）\par}
\typeout{TCELL:p335_1:5:\the\ht0:\the\dp0}
\setbox0=\vbox{\hsize=6.890cm\parindent=0pt\linespread{1}\fontsize{10}{13.50}\selectfont\raggedright\noindent 1.假药的判断标准，现采取实质标准，也即是否有疗效。有疗效，不算假药。以往采取形式标准，也即只要没有药监局批号，就是假药。\\{}2.假药的种类：\\{}（1）成分不符；\\{}（2）冒充的药品。以非药品冒充药品或者以他种药品冒充此种药品；\\{}（3）变质的药品；\\{}（4）标明的适应症超出规定范围。\\{}3.根据民间传统配方私自加工的药品，没有批号，但有疗效的，不算假药\par}
\typeout{TCELL:p335_1:6:\the\ht0:\the\dp0}
\setbox0=\vbox{\hsize=4.090cm\parindent=0pt\linespread{1}\fontsize{10}{13.50}\selectfont\raggedright\noindent 1.提供假药罪、提供劣药罪，是指药品使用单位或其人员明知是假药、劣药而提供给他人使用。提供，包括有偿提供和无偿提供。有偿提供就是销售。\\{}2.不以营利为目的，实施带有自救、互助性质的生产、进口、销售药品的行为，不认定为犯罪。这里的药品包括假药、劣药及没有批号的药。\\{}3.未取得药品销售许可证，销售药品，以前定非法经营罪，现在不定非法经营罪。如果符合其他罪的要件，定其他罪\par}
\typeout{TCELL:p335_1:7:\the\ht0:\the\dp0}
\setbox0=\vbox{\hsize=1.990cm\parindent=0pt\linespread{1}\fontsize{10}{13.50}\selectfont\centering\noindent 生产、销售、提供劣药罪\par}
\typeout{TCELL:p335_1:8:\the\ht0:\the\dp0}
\setbox0=\vbox{\hsize=1.990cm\parindent=0pt\linespread{1}\fontsize{10}{13.50}\selectfont\centering\noindent 实害犯\par}
\typeout{TCELL:p335_1:9:\the\ht0:\the\dp0}
\setbox0=\vbox{\hsize=6.890cm\parindent=0pt\linespread{1}\fontsize{10}{13.50}\selectfont\raggedright\noindent 劣药的种类：\\{}1.含量不符；\\{}2.被污染的药品；\\{}3.过期的药品；\\{}4.擅自添加防腐剂、辅料的药品\par}
\typeout{TCELL:p335_1:10:\the\ht0:\the\dp0}
\setbox0=\vbox{\hsize=1.990cm\parindent=0pt\linespread{1}\fontsize{10}{13.50}\selectfont\centering\noindent 妨害药品管理罪\par}
\typeout{TCELL:p335_1:11:\the\ht0:\the\dp0}
\setbox0=\vbox{\hsize=1.990cm\parindent=0pt\linespread{1}\fontsize{10}{13.50}\selectfont\centering\noindent 具体危险犯（足以严重危害人体健康）\par}
\typeout{TCELL:p335_1:12:\the\ht0:\the\dp0}
\setbox0=\vbox{\hsize=6.890cm\parindent=0pt\linespread{1}\fontsize{10}{13.50}\selectfont\raggedright\noindent 1.生产、销售国务院药品监督管理部门禁止使用的药品的；\\{}2.未取得药品相关批准证明文件生产、进口药品或者明知是上述药品而销售的；（若适应症、成分不明，或在境外未合法上市，视为有具体危险）\\{}3.药品申请注册中提供虚假的证明、数据、资料、样品或者采取其他欺骗手段的；\\{}4.编造生产、检验记录的（若影响药品的安全性、有效性，视为有具体危险）\par}
\typeout{TCELL:p335_1:13:\the\ht0:\the\dp0}
\setbox0=\vbox{\hsize=4.090cm\parindent=0pt\linespread{1}\fontsize{10}{13.50}\selectfont\raggedright\noindent 成立本罪，要求具体危险。\\{}私自进口、销售印度仿制药，如果没有具体危险，则不会构成该罪。因此，影片《我不是药神》中的男主角不会构成该罪，因为该药具有疗效\par}
\typeout{TCELL:p335_1:14:\the\ht0:\the\dp0}
\setbox0=\vbox{\hsize=1.990cm\parindent=0pt\linespread{1}\fontsize{10}{13.50}\selectfont\centering\noindent 药品监管渎职罪\par}
\typeout{TCELL:p335_1:15:\the\ht0:\the\dp0}
\setbox0=\vbox{\hsize=1.990cm\parindent=0pt\linespread{1}\fontsize{10}{13.50}\selectfont\centering\noindent 实害犯\par}
\typeout{TCELL:p335_1:16:\the\ht0:\the\dp0}
\setbox0=\vbox{\hsize=6.890cm\parindent=0pt\linespread{1}\fontsize{10}{13.50}\selectfont\raggedright\noindent 本罪既包括故意（滥用职权），也包括过失（玩忽职守）\par}
\typeout{TCELL:p335_1:17:\the\ht0:\the\dp0}
\setbox0=\vbox{\hsize=4.090cm\parindent=0pt\linespread{1}\fontsize{10}{13.50}\selectfont\raggedright\noindent 同时构成本罪和药品犯罪的共犯，想象竞合，择一重罪论处\par}
\typeout{TCELL:p335_1:18:\the\ht0:\the\dp0}
\setbox0=\vbox{\hsize=1.540cm\parindent=0pt\linespread{1}\fontsize{10}{13.50}\selectfont\bfseries\centering\noindent 物品种类\par}
\typeout{TCELL:p337_2:0:\the\ht0:\the\dp0}
\setbox0=\vbox{\hsize=13.940cm\parindent=0pt\linespread{1}\fontsize{10}{13.50}\selectfont\bfseries\centering\noindent 具体类型\par}
\typeout{TCELL:p337_2:1:\the\ht0:\the\dp0}
\setbox0=\vbox{\hsize=1.540cm\parindent=0pt\linespread{1}\fontsize{10}{13.50}\selectfont\centering\noindent 弹头、弹壳\par}
\typeout{TCELL:p337_2:2:\the\ht0:\the\dp0}
\setbox0=\vbox{\hsize=13.940cm\parindent=0pt\linespread{1}\fontsize{10}{13.50}\selectfont\raggedright\noindent 走私能够使用的弹头、弹壳，定走私弹药罪\par}
\typeout{TCELL:p337_2:3:\the\ht0:\the\dp0}
\setbox0=\vbox{\hsize=13.940cm\parindent=0pt\linespread{1}\fontsize{10}{13.50}\selectfont\raggedright\noindent 走私报废或者无法组装使用的弹头、弹壳，但不属于废物的，定走私普通货物、物品罪\par}
\typeout{TCELL:p337_2:4:\the\ht0:\the\dp0}
\setbox0=\vbox{\hsize=13.940cm\parindent=0pt\linespread{1}\fontsize{10}{13.50}\selectfont\raggedright\noindent 走私被鉴定为废物的弹头、弹壳，定走私废物罪\par}
\typeout{TCELL:p337_2:5:\the\ht0:\the\dp0}
\setbox0=\vbox{\hsize=1.540cm\parindent=0pt\linespread{1}\fontsize{10}{13.50}\selectfont\centering\noindent 仿真枪\par}
\typeout{TCELL:p337_2:6:\the\ht0:\the\dp0}
\setbox0=\vbox{\hsize=13.940cm\parindent=0pt\linespread{1}\fontsize{10}{13.50}\selectfont\raggedright\noindent 如果鉴定有杀伤力，达到真枪程度，则定走私武器罪\par}
\typeout{TCELL:p337_2:7:\the\ht0:\the\dp0}
\setbox0=\vbox{\hsize=13.940cm\parindent=0pt\linespread{1}\fontsize{10}{13.50}\selectfont\raggedright\noindent 如果鉴定有杀伤力，没达到真枪程度，但应当禁止进出口，则定走私国家禁止进出口的货物、物品罪\par}
\typeout{TCELL:p337_2:8:\the\ht0:\the\dp0}
\setbox0=\vbox{\hsize=13.940cm\parindent=0pt\linespread{1}\fontsize{10}{13.50}\selectfont\raggedright\noindent 如果鉴定没有杀伤力，允许进出口，则定走私普通货物、物品罪\par}
\typeout{TCELL:p337_2:9:\the\ht0:\the\dp0}
\setbox0=\vbox{\hsize=2.890cm\parindent=0pt\linespread{1}\fontsize{10}{13.50}\selectfont\bfseries\centering\noindent 罪名\par}
\typeout{TCELL:p338_1:0:\the\ht0:\the\dp0}
\setbox0=\vbox{\hsize=5.690cm\parindent=0pt\linespread{1}\fontsize{10}{13.50}\selectfont\bfseries\centering\noindent 受贿罪\par}
\typeout{TCELL:p338_1:1:\the\ht0:\the\dp0}
\setbox0=\vbox{\hsize=6.640cm\parindent=0pt\linespread{1}\fontsize{10}{13.50}\selectfont\bfseries\centering\noindent 非国家工作人员受贿罪\par}
\typeout{TCELL:p338_1:2:\the\ht0:\the\dp0}
\setbox0=\vbox{\hsize=2.890cm\parindent=0pt\linespread{1}\fontsize{10}{13.50}\selectfont\centering\noindent 主体身份\par}
\typeout{TCELL:p338_1:3:\the\ht0:\the\dp0}
\setbox0=\vbox{\hsize=5.690cm\parindent=0pt\linespread{1}\fontsize{10}{13.50}\selectfont\centering\noindent 国家工作人员\par}
\typeout{TCELL:p338_1:4:\the\ht0:\the\dp0}
\setbox0=\vbox{\hsize=6.640cm\parindent=0pt\linespread{1}\fontsize{10}{13.50}\selectfont\centering\noindent 普通单位的工作人员（非国家工作人员）\par}
\typeout{TCELL:p338_1:5:\the\ht0:\the\dp0}
\setbox0=\vbox{\hsize=2.890cm\parindent=0pt\linespread{1}\fontsize{10}{13.50}\selectfont\centering\noindent 行为方式\par}
\typeout{TCELL:p338_1:6:\the\ht0:\the\dp0}
\setbox0=\vbox{\hsize=5.690cm\parindent=0pt\linespread{1}\fontsize{10}{13.50}\selectfont\raggedright\noindent 索取贿赂不要求为他人谋取利益。\\{}只有收受贿赂要求为他人谋取利益\par}
\typeout{TCELL:p338_1:7:\the\ht0:\the\dp0}
\setbox0=\vbox{\hsize=6.640cm\parindent=0pt\linespread{1}\fontsize{10}{13.50}\selectfont\raggedright\noindent 不管是索取贿赂，还是收受贿赂，都要求为他人谋取利益\par}
\typeout{TCELL:p338_1:8:\the\ht0:\the\dp0}
\setbox0=\vbox{\hsize=2.890cm\parindent=0pt\linespread{1}\fontsize{10}{13.50}\selectfont\centering\noindent 为他人谋取利益\par}
\typeout{TCELL:p338_1:9:\the\ht0:\the\dp0}
\setbox0=\vbox{\hsize=12.590cm\parindent=0pt\linespread{1}\fontsize{10}{13.50}\selectfont\raggedright\noindent “为他人谋取利益”，既包括正当利益，也包括不正当利益。“为他人谋取利益”只要求许诺为他人谋取利益，不要求实际上为他人谋取了利益\par}
\typeout{TCELL:p338_1:10:\the\ht0:\the\dp0}
\setbox0=\vbox{\hsize=1.990cm\parindent=0pt\linespread{1}\fontsize{10}{13.50}\selectfont\bfseries\centering\noindent 罪名区分\par}
\typeout{TCELL:p339_1:0:\the\ht0:\the\dp0}
\setbox0=\vbox{\hsize=7.440cm\parindent=0pt\linespread{1}\fontsize{10}{13.50}\selectfont\bfseries\centering\noindent 区分标准\par}
\typeout{TCELL:p339_1:1:\the\ht0:\the\dp0}
\setbox0=\vbox{\hsize=5.790cm\parindent=0pt\linespread{1}\fontsize{10}{13.50}\selectfont\bfseries\centering\noindent 要点\par}
\typeout{TCELL:p339_1:2:\the\ht0:\the\dp0}
\setbox0=\vbox{\hsize=1.990cm\parindent=0pt\linespread{1}\fontsize{10}{13.50}\selectfont\centering\noindent 伪造货币罪\par}
\typeout{TCELL:p339_1:3:\the\ht0:\the\dp0}
\setbox0=\vbox{\hsize=7.440cm\parindent=0pt\linespread{1}\fontsize{10}{13.50}\selectfont\raggedright\noindent 方式1：用假材料制作假币\par}
\typeout{TCELL:p339_1:4:\the\ht0:\the\dp0}
\setbox0=\vbox{\hsize=7.440cm\parindent=0pt\linespread{1}\fontsize{10}{13.50}\selectfont\raggedright\noindent 方式2：用真币材料制作假币，并且对真币材料完全翻新，不保留原真币的基础。\\{}例1，将真币的碎片作为材料，拼凑成新的假币，属于伪造货币。\\{}例2，将金属货币熔化后，制作成新的假币，属于伪造货币\par}
\typeout{TCELL:p339_1:5:\the\ht0:\the\dp0}
\setbox0=\vbox{\hsize=5.790cm\parindent=0pt\linespread{1}\fontsize{10}{13.50}\selectfont\raggedright\noindent 1.只包括正在流通的货币。伪造停止流通的货币并使用的，以诈骗罪论处。（2011年第59题）\\{}2.包括正在流通的境外货币\par}
\typeout{TCELL:p339_1:6:\the\ht0:\the\dp0}
\setbox0=\vbox{\hsize=1.990cm\parindent=0pt\linespread{1}\fontsize{10}{13.50}\selectfont\centering\noindent 变造货币罪\par}
\typeout{TCELL:p339_1:7:\the\ht0:\the\dp0}
\setbox0=\vbox{\hsize=7.440cm\parindent=0pt\linespread{1}\fontsize{10}{13.50}\selectfont\raggedright\noindent 用真币材料制作假币，并且保留原真币的基础。\\{}例1，将硬币周边的金属刮下来，减少硬币金属含量，属于变造。例2，将100元面额的人民币数字改为50元，属于变造。（2013年第14题）\\{}[注意］同时采用伪造和变造的手段，拼凑假币的行为，以伪造货币来论\par}
\typeout{TCELL:p339_1:8:\the\ht0:\the\dp0}
\setbox0=\vbox{\hsize=5.790cm\parindent=0pt\linespread{1}\fontsize{10}{13.50}\selectfont\raggedright\noindent “出售、购买、运输假币罪”“持有、使用假币罪”中的假币，均指伪造的货币，不包括变造的货币（2016年第54题）\par}
\typeout{TCELL:p339_1:9:\the\ht0:\the\dp0}
\setbox0=\vbox{\hsize=8.090cm\parindent=0pt\linespread{1}\fontsize{10}{13.50}\selectfont\bfseries\centering\noindent 同一对象\par}
\typeout{TCELL:p341_1:0:\the\ht0:\the\dp0}
\setbox0=\vbox{\hsize=7.390cm\parindent=0pt\linespread{1}\fontsize{10}{13.50}\selectfont\bfseries\centering\noindent 不同对象\par}
\typeout{TCELL:p341_1:1:\the\ht0:\the\dp0}
\setbox0=\vbox{\hsize=8.090cm\parindent=0pt\linespread{1}\fontsize{10}{13.50}\selectfont\raggedright\noindent 伪造+出售、购买、运输、持有、使用=伪造货币罪\\{}（这是“先干什么、后干什么”的串联关系，前罪吸收后罪）\par}
\typeout{TCELL:p341_1:2:\the\ht0:\the\dp0}
\setbox0=\vbox{\hsize=7.390cm\parindent=0pt\linespread{1}\fontsize{10}{13.50}\selectfont\raggedright\noindent 伪造A假币+出售、购买、运输、持有、使用B假币，并罚\\{}（前后两个罪是独立罪名，所以并罚）（2013年第14题）\par}
\typeout{TCELL:p341_1:3:\the\ht0:\the\dp0}
\setbox0=\vbox{\hsize=8.090cm\parindent=0pt\linespread{1}\fontsize{10}{13.50}\selectfont\raggedright\noindent 出售+购买+运输，定出售、购买、运输假币罪\par}
\typeout{TCELL:p341_1:4:\the\ht0:\the\dp0}
\setbox0=\vbox{\hsize=7.390cm\parindent=0pt\linespread{1}\fontsize{10}{13.50}\selectfont\raggedright\noindent 出售A假币+购买B假币+运输C假币，定出售、购买、运输假币罪，不并罚（这是选择性罪名，无论对象是否同一，都不并罚）\par}
\typeout{TCELL:p341_1:5:\the\ht0:\the\dp0}
\setbox0=\vbox{\hsize=8.090cm\parindent=0pt\linespread{1}\fontsize{10}{13.50}\selectfont\raggedright\noindent 出售/运输+使用，并罚（这是“边干什么、边干什么”的并联关系，所以并罚）\par}
\typeout{TCELL:p341_1:6:\the\ht0:\the\dp0}
\setbox0=\vbox{\hsize=7.390cm\parindent=0pt\linespread{1}\fontsize{10}{13.50}\selectfont\raggedright\noindent 出售/运输A假币+使用B假币，并罚（前后两个罪是独立罪名，所以并罚）（2010年第13题）\par}
\typeout{TCELL:p341_1:7:\the\ht0:\the\dp0}
\setbox0=\vbox{\hsize=8.090cm\parindent=0pt\linespread{1}\fontsize{10}{13.50}\selectfont\raggedright\noindent 购买+使用=购买假币罪（这是串联关系，前罪吸收后罪)\par}
\typeout{TCELL:p341_1:8:\the\ht0:\the\dp0}
\setbox0=\vbox{\hsize=7.390cm\parindent=0pt\linespread{1}\fontsize{10}{13.50}\selectfont\raggedright\noindent 购买A假币+使用B假币，并罚（前后两个罪是独立罪名，所以并罚）\par}
\typeout{TCELL:p341_1:9:\the\ht0:\the\dp0}
\setbox0=\vbox{\hsize=1.740cm\parindent=0pt\linespread{1}\fontsize{10}{13.50}\selectfont\bfseries\centering\noindent 构成要件\par}
\typeout{TCELL:p343_1:0:\the\ht0:\the\dp0}
\setbox0=\vbox{\hsize=4.140cm\parindent=0pt\linespread{1}\fontsize{10}{13.50}\selectfont\bfseries\centering\noindent 具体内容\par}
\typeout{TCELL:p343_1:1:\the\ht0:\the\dp0}
\setbox0=\vbox{\hsize=9.340cm\parindent=0pt\linespread{1}\fontsize{10}{13.50}\selectfont\bfseries\centering\noindent 要点\par}
\typeout{TCELL:p343_1:2:\the\ht0:\the\dp0}
\setbox0=\vbox{\hsize=1.740cm\parindent=0pt\linespread{1}\fontsize{10}{13.50}\selectfont\centering\noindent 行为主体\par}
\typeout{TCELL:p343_1:3:\the\ht0:\the\dp0}
\setbox0=\vbox{\hsize=4.140cm\parindent=0pt\linespread{1}\fontsize{10}{13.50}\selectfont\raggedright\noindent 不限于贷款人，包括担保人\par}
\typeout{TCELL:p343_1:4:\the\ht0:\the\dp0}
\setbox0=\vbox{\hsize=9.340cm\parindent=0pt\linespread{1}\fontsize{10}{13.50}\selectfont\raggedright\noindent 甲没有非法占有目的，为了取得银行贷款，请乙提供担保。乙提供了虚假的担保材料，甲不知情，银行因为受骗而发放贷款。甲因客观原因无法还款，银行遭受重大损失。乙独立构成本罪\par}
\typeout{TCELL:p343_1:5:\the\ht0:\the\dp0}
\setbox0=\vbox{\hsize=1.740cm\parindent=0pt\linespread{1}\fontsize{10}{13.50}\selectfont\centering\noindent 行为方式\par}
\typeout{TCELL:p343_1:6:\the\ht0:\the\dp0}
\setbox0=\vbox{\hsize=4.140cm\parindent=0pt\linespread{1}\fontsize{10}{13.50}\selectfont\raggedright\noindent 以欺骗手段取得银行贷款\par}
\typeout{TCELL:p343_1:7:\the\ht0:\the\dp0}
\setbox0=\vbox{\hsize=9.340cm\parindent=0pt\linespread{1}\fontsize{10}{13.50}\selectfont\raggedright\noindent 如果金融机构工作人员知道真相仍然放贷，则行为人不构成本罪，相反，金融机构工作人员应构成违法发放贷款罪\par}
\typeout{TCELL:p343_1:8:\the\ht0:\the\dp0}
\setbox0=\vbox{\hsize=1.740cm\parindent=0pt\linespread{1}\fontsize{10}{13.50}\selectfont\centering\noindent 危害结果\par}
\typeout{TCELL:p343_1:9:\the\ht0:\the\dp0}
\setbox0=\vbox{\hsize=4.140cm\parindent=0pt\linespread{1}\fontsize{10}{13.50}\selectfont\raggedright\noindent 要求给银行造成重大损失\par}
\typeout{TCELL:p343_1:10:\the\ht0:\the\dp0}
\setbox0=\vbox{\hsize=9.340cm\parindent=0pt\linespread{1}\fontsize{10}{13.50}\selectfont\raggedright\noindent 该罪的成立条件原来包括两种情形，或者给银行造成重大损失，或者有严重情节。《刑法修正案（十一）》删除“有严重情节”。本罪的成立条件只剩下“给银行造成重大损失”\par}
\typeout{TCELL:p343_1:11:\the\ht0:\the\dp0}
\setbox0=\vbox{\hsize=1.740cm\parindent=0pt\linespread{1}\fontsize{10}{13.50}\selectfont\centering\noindent 主观要件\par}
\typeout{TCELL:p343_1:12:\the\ht0:\the\dp0}
\setbox0=\vbox{\hsize=4.140cm\parindent=0pt\linespread{1}\fontsize{10}{13.50}\selectfont\raggedright\noindent 对贷款没有非法占有目的，也即具有归还意思\par}
\typeout{TCELL:p343_1:13:\the\ht0:\the\dp0}
\setbox0=\vbox{\hsize=9.340cm\parindent=0pt\linespread{1}\fontsize{10}{13.50}\selectfont\raggedright\noindent 本罪与贷款诈骗罪的区别在于：贷款诈骗罪具有非法占有目的\par}
\typeout{TCELL:p343_1:14:\the\ht0:\the\dp0}
\setbox0=\vbox{\hsize=5.790cm\parindent=0pt\linespread{1}\fontsize{10}{13.50}\selectfont\bfseries\centering\noindent 非法取得行为\par}
\typeout{TCELL:p349_1:0:\the\ht0:\the\dp0}
\setbox0=\vbox{\hsize=4.140cm\parindent=0pt\linespread{1}\fontsize{10}{13.50}\selectfont\bfseries\centering\noindent 使用行为\par}
\typeout{TCELL:p349_1:1:\the\ht0:\the\dp0}
\setbox0=\vbox{\hsize=5.290cm\parindent=0pt\linespread{1}\fontsize{10}{13.50}\selectfont\bfseries\centering\noindent 结论\par}
\typeout{TCELL:p349_1:2:\the\ht0:\the\dp0}
\setbox0=\vbox{\hsize=5.790cm\parindent=0pt\linespread{1}\fontsize{10}{13.50}\selectfont\centering\noindent 伪造信用卡（伪造金融票证罪）\par}
\typeout{TCELL:p349_1:3:\the\ht0:\the\dp0}
\setbox0=\vbox{\hsize=4.140cm\parindent=0pt\linespread{1}\fontsize{10}{13.50}\selectfont\centering\noindent 无论对人或机器使用\par}
\typeout{TCELL:p349_1:4:\the\ht0:\the\dp0}
\setbox0=\vbox{\hsize=5.290cm\parindent=0pt\linespread{1}\fontsize{10}{13.50}\selectfont\centering\noindent 信用卡诈骗罪（牵连犯）\par}
\typeout{TCELL:p349_1:5:\the\ht0:\the\dp0}
\setbox0=\vbox{\hsize=5.790cm\parindent=0pt\linespread{1}\fontsize{10}{13.50}\selectfont\centering\noindent 以虚假的身份证明骗领信用卡（妨害信用卡管理罪）\par}
\typeout{TCELL:p349_1:6:\the\ht0:\the\dp0}
\setbox0=\vbox{\hsize=5.290cm\parindent=0pt\linespread{1}\fontsize{10}{13.50}\selectfont\centering\noindent 信用卡诈骗罪（牵连犯）\par}
\typeout{TCELL:p349_1:7:\the\ht0:\the\dp0}
\setbox0=\vbox{\hsize=5.790cm\parindent=0pt\linespread{1}\fontsize{10}{13.50}\selectfont\bfseries\centering\noindent 非法取得行为\par}
\typeout{TCELL:p349_2:0:\the\ht0:\the\dp0}
\setbox0=\vbox{\hsize=4.140cm\parindent=0pt\linespread{1}\fontsize{10}{13.50}\selectfont\bfseries\centering\noindent 使用行为\par}
\typeout{TCELL:p349_2:1:\the\ht0:\the\dp0}
\setbox0=\vbox{\hsize=5.290cm\parindent=0pt\linespread{1}\fontsize{10}{13.50}\selectfont\bfseries\centering\noindent 结论\par}
\typeout{TCELL:p349_2:2:\the\ht0:\the\dp0}
\setbox0=\vbox{\hsize=5.790cm\parindent=0pt\linespread{1}\fontsize{10}{13.50}\selectfont\centering\noindent 侵占，诈骗，抢夺、敲诈\par}
\typeout{TCELL:p349_2:3:\the\ht0:\the\dp0}
\setbox0=\vbox{\hsize=4.140cm\parindent=0pt\linespread{1}\fontsize{10}{13.50}\selectfont\centering\noindent 无论对人或机器使用\par}
\typeout{TCELL:p349_2:4:\the\ht0:\the\dp0}
\setbox0=\vbox{\hsize=5.290cm\parindent=0pt\linespread{1}\fontsize{10}{13.50}\selectfont\centering\noindent 信用卡诈骗罪\par}
\typeout{TCELL:p349_2:5:\the\ht0:\the\dp0}
\setbox0=\vbox{\hsize=5.790cm\parindent=0pt\linespread{1}\fontsize{10}{13.50}\selectfont\centering\noindent 抢劫\par}
\typeout{TCELL:p349_2:6:\the\ht0:\the\dp0}
\setbox0=\vbox{\hsize=5.290cm\parindent=0pt\linespread{1}\fontsize{10}{13.50}\selectfont\centering\noindent 抢劫罪\par}
\typeout{TCELL:p349_2:7:\the\ht0:\the\dp0}
\setbox0=\vbox{\hsize=5.790cm\parindent=0pt\linespread{1}\fontsize{10}{13.50}\selectfont\centering\noindent 盗窃\par}
\typeout{TCELL:p349_2:8:\the\ht0:\the\dp0}
\setbox0=\vbox{\hsize=5.290cm\parindent=0pt\linespread{1}\fontsize{10}{13.50}\selectfont\centering\noindent 盗窃罪\par}
\typeout{TCELL:p349_2:9:\the\ht0:\the\dp0}
\setbox0=\vbox{\hsize=6.590cm\parindent=0pt\linespread{1}\fontsize{10}{13.50}\selectfont\bfseries\centering\noindent 非法取得行为\par}
\typeout{TCELL:p350_1:0:\the\ht0:\the\dp0}
\setbox0=\vbox{\hsize=4.690cm\parindent=0pt\linespread{1}\fontsize{10}{13.50}\selectfont\bfseries\centering\noindent 使用行为\par}
\typeout{TCELL:p350_1:1:\the\ht0:\the\dp0}
\setbox0=\vbox{\hsize=3.940cm\parindent=0pt\linespread{1}\fontsize{10}{13.50}\selectfont\bfseries\centering\noindent 结论\par}
\typeout{TCELL:p350_1:2:\the\ht0:\the\dp0}
\setbox0=\vbox{\hsize=6.590cm\parindent=0pt\linespread{1}\fontsize{10}{13.50}\selectfont\centering\noindent 一切非法方式（窃取、收买信用卡信息罪）\par}
\typeout{TCELL:p350_1:3:\the\ht0:\the\dp0}
\setbox0=\vbox{\hsize=4.690cm\parindent=0pt\linespread{1}\fontsize{10}{13.50}\selectfont\raggedright\noindent 无论对人或机器（网络）使用\par}
\typeout{TCELL:p350_1:4:\the\ht0:\the\dp0}
\setbox0=\vbox{\hsize=3.940cm\parindent=0pt\linespread{1}\fontsize{10}{13.50}\selectfont\raggedright\noindent 信用卡诈骗罪（牵连犯）\par}
\typeout{TCELL:p350_1:5:\the\ht0:\the\dp0}
\setbox0=\vbox{\hsize=8.590cm\parindent=0pt\linespread{1}\fontsize{10}{13.50}\selectfont\bfseries\centering\noindent 类型\par}
\typeout{TCELL:p350_2:0:\the\ht0:\the\dp0}
\setbox0=\vbox{\hsize=6.890cm\parindent=0pt\linespread{1}\fontsize{10}{13.50}\selectfont\bfseries\centering\noindent 结论\par}
\typeout{TCELL:p350_2:1:\the\ht0:\the\dp0}
\setbox0=\vbox{\hsize=8.590cm\parindent=0pt\linespread{1}\fontsize{10}{13.50}\selectfont\centering\noindent 盗用微信、支付宝信息资料\par}
\typeout{TCELL:p350_2:2:\the\ht0:\the\dp0}
\setbox0=\vbox{\hsize=6.890cm\parindent=0pt\linespread{1}\fontsize{10}{13.50}\selectfont\raggedright\noindent 盗窃罪\par}
\typeout{TCELL:p350_2:3:\the\ht0:\the\dp0}
\setbox0=\vbox{\hsize=8.590cm\parindent=0pt\linespread{1}\fontsize{10}{13.50}\selectfont\centering\noindent 冒用蚂蚁花呗、借呗\par}
\typeout{TCELL:p350_2:4:\the\ht0:\the\dp0}
\setbox0=\vbox{\hsize=6.890cm\parindent=0pt\linespread{1}\fontsize{10}{13.50}\selectfont\raggedright\noindent 贷款诈骗罪\par}
\typeout{TCELL:p350_2:5:\the\ht0:\the\dp0}
\setbox0=\vbox{\hsize=1.990cm\parindent=0pt\linespread{1}\fontsize{10}{13.50}\selectfont\bfseries\centering\noindent 保险诈骗罪\par}
\typeout{TCELL:p351_1:0:\the\ht0:\the\dp0}
\setbox0=\vbox{\hsize=3.940cm\parindent=0pt\linespread{1}\fontsize{10}{13.50}\selectfont\bfseries\centering\noindent 内容\par}
\typeout{TCELL:p351_1:1:\the\ht0:\the\dp0}
\setbox0=\vbox{\hsize=9.290cm\parindent=0pt\linespread{1}\fontsize{10}{13.50}\selectfont\bfseries\centering\noindent 范例\par}
\typeout{TCELL:p351_1:2:\the\ht0:\the\dp0}
\setbox0=\vbox{\hsize=1.990cm\parindent=0pt\linespread{1}\fontsize{10}{13.50}\selectfont\centering\noindent 行为主体\par}
\typeout{TCELL:p351_1:3:\the\ht0:\the\dp0}
\setbox0=\vbox{\hsize=3.940cm\parindent=0pt\linespread{1}\fontsize{10}{13.50}\selectfont\raggedright\noindent 是特殊主体，限于投保人、被保险人、受益人。不具有主体身份的人骗取保险金的，可构成普通的诈骗罪\par}
\typeout{TCELL:p351_1:4:\the\ht0:\the\dp0}
\setbox0=\vbox{\hsize=9.290cm\parindent=0pt\linespread{1}\fontsize{10}{13.50}\selectfont\raggedright\noindent 例1，个体户甲开办的汽车修理厂系某保险公司指定的汽车修理厂家。甲在为他人修理汽车时，多次夸大汽车毁损程度，向保险公司多报汽车修理费用，从保险公司骗取12万余元。对甲的行为应以诈骗罪论处。（2004年第5题）\\{}例2，甲公司给经理王某配了一辆轿车，轿车登记在公司名下，由王某使用。王某将车私下变卖，对甲公司领导谎称被盗，让甲公司去申请理赔。甲公司不知情，到保险公司理赔，获得5万元。王某构成职务侵占罪。王某对保险公司不构成保险诈骗罪的间接正犯，因为王某不具有投保人的身份，但可以构成诈骗罪的间接正犯\par}
\typeout{TCELL:p351_1:5:\the\ht0:\the\dp0}
\setbox0=\vbox{\hsize=3.940cm\parindent=0pt\linespread{1}\fontsize{10}{13.50}\selectfont\centering\noindent 包括单位\par}
\typeout{TCELL:p351_1:6:\the\ht0:\the\dp0}
\setbox0=\vbox{\hsize=9.290cm\parindent=0pt\linespread{1}\fontsize{10}{13.50}\selectfont\raggedright\noindent 某公司决定诈骗保险金，故意放火烧毁公司的仓库，并危害了公共安全，然后骗取保险金。对该公司定保险诈骗罪。由于单位不能构成放火罪，对公司直接责任人员定放火罪\par}
\typeout{TCELL:p351_1:7:\the\ht0:\the\dp0}
\setbox0=\vbox{\hsize=1.990cm\parindent=0pt\linespread{1}\fontsize{10}{13.50}\selectfont\centering\noindent 侵害对象\par}
\typeout{TCELL:p351_1:8:\the\ht0:\the\dp0}
\setbox0=\vbox{\hsize=3.940cm\parindent=0pt\linespread{1}\fontsize{10}{13.50}\selectfont\raggedright\noindent 商业保险。不包括社会保险\par}
\typeout{TCELL:p351_1:9:\the\ht0:\the\dp0}
\setbox0=\vbox{\hsize=9.290cm\parindent=0pt\linespread{1}\fontsize{10}{13.50}\selectfont\raggedright\noindent 甲拾得乙的社保卡。甲在医院收费处交费时，冒用乙的社保卡。甲构成诈骗罪，不构成保险诈骗罪\par}
\typeout{TCELL:p351_1:10:\the\ht0:\the\dp0}
\setbox0=\vbox{\hsize=1.990cm\parindent=0pt\linespread{1}\fontsize{10}{13.50}\selectfont\centering\noindent 着手的标准\par}
\typeout{TCELL:p351_1:11:\the\ht0:\the\dp0}
\setbox0=\vbox{\hsize=13.490cm\parindent=0pt\linespread{1}\fontsize{10}{13.50}\selectfont\raggedright\noindent 虚构保险标的、制造保险事故，属于预备行为；只有到保险公司理赔，才是着手，才是实行行为\par}
\typeout{TCELL:p351_1:12:\the\ht0:\the\dp0}
\setbox0=\vbox{\hsize=1.990cm\parindent=0pt\linespread{1}\fontsize{10}{13.50}\selectfont\bfseries\centering\noindent 保险诈骗罪\par}
\typeout{TCELL:p352_1:0:\the\ht0:\the\dp0}
\setbox0=\vbox{\hsize=3.940cm\parindent=0pt\linespread{1}\fontsize{10}{13.50}\selectfont\bfseries\centering\noindent 内容\par}
\typeout{TCELL:p352_1:1:\the\ht0:\the\dp0}
\setbox0=\vbox{\hsize=9.290cm\parindent=0pt\linespread{1}\fontsize{10}{13.50}\selectfont\bfseries\centering\noindent 范例\par}
\typeout{TCELL:p352_1:2:\the\ht0:\the\dp0}
\setbox0=\vbox{\hsize=1.990cm\parindent=0pt\linespread{1}\fontsize{10}{13.50}\selectfont\centering\noindent 罪数问题\par}
\typeout{TCELL:p352_1:3:\the\ht0:\the\dp0}
\setbox0=\vbox{\hsize=3.940cm\parindent=0pt\linespread{1}\fontsize{10}{13.50}\selectfont\centering\noindent 前后两个行为\par}
\typeout{TCELL:p352_1:4:\the\ht0:\the\dp0}
\setbox0=\vbox{\hsize=9.290cm\parindent=0pt\linespread{1}\fontsize{10}{13.50}\selectfont\raggedright\noindent 甲将妻子乙作为被保险人进行投保，然后故意杀害乙，并骗取保险金，构成故意杀人罪和保险诈骗罪，数罪并罚①\par}
\typeout{TCELL:p352_1:5:\the\ht0:\the\dp0}
\setbox0=\vbox{\hsize=3.940cm\parindent=0pt\linespread{1}\fontsize{10}{13.50}\selectfont\centering\noindent 一个行为\par}
\typeout{TCELL:p352_1:6:\the\ht0:\the\dp0}
\setbox0=\vbox{\hsize=9.290cm\parindent=0pt\linespread{1}\fontsize{10}{13.50}\selectfont\raggedright\noindent 甲将妻子乙作为被保险人进行投保，然后杀死乙，尚未理赔时被抓。甲既构成故意杀人罪（既遂），又构成保险诈骗罪（预备），属于想象竞合，择一重罪论处\par}
\typeout{TCELL:p352_1:7:\the\ht0:\the\dp0}
\setbox0=\vbox{\hsize=2.540cm\parindent=0pt\linespread{1}\fontsize{10}{13.50}\selectfont\bfseries\centering\noindent 罪名\par}
\typeout{TCELL:p352_2:0:\the\ht0:\the\dp0}
\setbox0=\vbox{\hsize=6.690cm\parindent=0pt\linespread{1}\fontsize{10}{13.50}\selectfont\bfseries\centering\noindent 构成要件\par}
\typeout{TCELL:p352_2:1:\the\ht0:\the\dp0}
\setbox0=\vbox{\hsize=5.990cm\parindent=0pt\linespread{1}\fontsize{10}{13.50}\selectfont\bfseries\centering\noindent 要点\par}
\typeout{TCELL:p352_2:2:\the\ht0:\the\dp0}
\setbox0=\vbox{\hsize=2.540cm\parindent=0pt\linespread{1}\fontsize{10}{13.50}\selectfont\centering\noindent 非法吸收公众存款罪\par}
\typeout{TCELL:p352_2:3:\the\ht0:\the\dp0}
\setbox0=\vbox{\hsize=6.690cm\parindent=0pt\linespread{1}\fontsize{10}{13.50}\selectfont\raggedright\noindent 1.非法性。这是指未经批准。\\{}2.公开性。这是指通过媒体等途径向社会公开宣传。\\{}3.利诱性。这是指承诺在一定期限内还本付息或者给付回报。\\{}4.公众性。这是指向社会公众即社会不特定对象吸收资金。\\{}5.存款性。这是指，行为人吸收存款是想用于放贷等金融活动，由此吸收的资金才能称为吸收存款。如果吸收资金用于正当的生产经营活动，则仅属于吸收资金，不属于吸收存款\par}
\typeout{TCELL:p352_2:4:\the\ht0:\the\dp0}
\setbox0=\vbox{\hsize=5.990cm\parindent=0pt\linespread{1}\fontsize{10}{13.50}\selectfont\raggedright\noindent 1.向亲友或者单位内部人员吸收资金的，不属于向社会公众吸收存款；但如果明知亲友或者单位内部人员又向不特定对象吸收资金而予以放任的，属于向社会公众吸收存款。（2014年第14题）\\{}2.通过传销手段实施，同时触犯组织、领导传销活动罪和非法吸收公众存款罪或集资诈骗罪的，属于牵连犯，择一重罪论处\par}
\typeout{TCELL:p352_2:5:\the\ht0:\the\dp0}
\setbox0=\vbox{\hsize=2.540cm\parindent=0pt\linespread{1}\fontsize{10}{13.50}\selectfont\centering\noindent 集资诈骗罪\par}
\typeout{TCELL:p352_2:6:\the\ht0:\the\dp0}
\setbox0=\vbox{\hsize=6.690cm\parindent=0pt\linespread{1}\fontsize{10}{13.50}\selectfont\raggedright\noindent 通过非法集资的方式实施的诈骗罪\par}
\typeout{TCELL:p352_2:7:\the\ht0:\the\dp0}
\setbox0=\vbox{\hsize=5.990cm\parindent=0pt\linespread{1}\fontsize{10}{13.50}\selectfont\raggedright\noindent 1.集资诈骗罪有非法占有目的。非法吸收公众存款罪没有非法占有目的。\\{}2.集资后不用于生产经营，或携带集资款逃匿，视为具有非法占有目的\par}
\typeout{TCELL:p352_2:8:\the\ht0:\the\dp0}
\setbox0=\vbox{\hsize=2.540cm\parindent=0pt\linespread{1}\fontsize{10}{13.50}\selectfont\centering\noindent 擅自发行股票、公司、企业债券罪\par}
\typeout{TCELL:p352_2:9:\the\ht0:\the\dp0}
\setbox0=\vbox{\hsize=6.690cm\parindent=0pt\linespread{1}\fontsize{10}{13.50}\selectfont\raggedright\noindent 未经批准，向社会不特定对象发行、以转让股权等方式变相发行股票或者公司、企业债券\par}
\typeout{TCELL:p352_2:10:\the\ht0:\the\dp0}
\setbox0=\vbox{\hsize=5.990cm\parindent=0pt\linespread{1}\fontsize{10}{13.50}\selectfont\raggedright\noindent 针对特定对象发行，特定对象人数超过200人的，也构成本罪\par}
\typeout{TCELL:p352_2:11:\the\ht0:\the\dp0}
\setbox0=\vbox{\hsize=2.540cm\parindent=0pt\linespread{1}\fontsize{10}{13.50}\selectfont\bfseries\centering\noindent 罪名\par}
\typeout{TCELL:p353_1:0:\the\ht0:\the\dp0}
\setbox0=\vbox{\hsize=6.690cm\parindent=0pt\linespread{1}\fontsize{10}{13.50}\selectfont\bfseries\centering\noindent 构成要件\par}
\typeout{TCELL:p353_1:1:\the\ht0:\the\dp0}
\setbox0=\vbox{\hsize=5.990cm\parindent=0pt\linespread{1}\fontsize{10}{13.50}\selectfont\bfseries\centering\noindent 要点\par}
\typeout{TCELL:p353_1:2:\the\ht0:\the\dp0}
\setbox0=\vbox{\hsize=2.540cm\parindent=0pt\linespread{1}\fontsize{10}{13.50}\selectfont\centering\noindent 非法经营罪\par}
\typeout{TCELL:p353_1:3:\the\ht0:\the\dp0}
\setbox0=\vbox{\hsize=6.690cm\parindent=0pt\linespread{1}\fontsize{10}{13.50}\selectfont\raggedright\noindent 1.未经批准，非法经营证券、期货、保险业务。\\{}2.未经核准，擅自发行基金。\\{}3.中介机构非法代理买卖非上市公司股票\par}
\typeout{TCELL:p353_1:4:\the\ht0:\the\dp0}
\setbox0=\vbox{\hsize=5.990cm\parindent=0pt\linespread{1}\fontsize{10}{13.50}\selectfont\raggedright\noindent 成立本罪，同时构成其他犯罪，想象竞合，择一重罪论处\par}
\typeout{TCELL:p353_1:5:\the\ht0:\the\dp0}
\setbox0=\vbox{\hsize=1.140cm\parindent=0pt\linespread{1}\fontsize{10}{13.50}\selectfont\bfseries\centering\noindent 传播\par}
\typeout{TCELL:p355_1:0:\the\ht0:\the\dp0}
\setbox0=\vbox{\hsize=1.490cm\parindent=0pt\linespread{1}\fontsize{10}{13.50}\selectfont\bfseries\centering\noindent 方式\par}
\typeout{TCELL:p355_1:1:\the\ht0:\the\dp0}
\setbox0=\vbox{\hsize=4.940cm\parindent=0pt\linespread{1}\fontsize{10}{13.50}\selectfont\bfseries\centering\noindent 含义\par}
\typeout{TCELL:p355_1:2:\the\ht0:\the\dp0}
\setbox0=\vbox{\hsize=3.540cm\parindent=0pt\linespread{1}\fontsize{10}{13.50}\selectfont\bfseries\centering\noindent 线上\par}
\typeout{TCELL:p355_1:3:\the\ht0:\the\dp0}
\setbox0=\vbox{\hsize=3.590cm\parindent=0pt\linespread{1}\fontsize{10}{13.50}\selectfont\bfseries\centering\noindent 线下\par}
\typeout{TCELL:p355_1:4:\the\ht0:\the\dp0}
\setbox0=\vbox{\hsize=1.490cm\parindent=0pt\linespread{1}\fontsize{10}{13.50}\selectfont\centering\noindent 发行\par}
\typeout{TCELL:p355_1:5:\the\ht0:\the\dp0}
\setbox0=\vbox{\hsize=4.940cm\parindent=0pt\linespread{1}\fontsize{10}{13.50}\selectfont\raggedright\noindent 将侵权作品（内容+载体）转让给不特定人或多数人\par}
\typeout{TCELL:p355_1:6:\the\ht0:\the\dp0}
\setbox0=\vbox{\hsize=3.540cm\parindent=0pt\linespread{1}\fontsize{10}{13.50}\selectfont\raggedright\noindent 线上发行，构成本罪（例1）\par}
\typeout{TCELL:p355_1:7:\the\ht0:\the\dp0}
\setbox0=\vbox{\hsize=3.590cm\parindent=0pt\linespread{1}\fontsize{10}{13.50}\selectfont\raggedright\noindent 线下发行，构成本罪（例2）\par}
\typeout{TCELL:p355_1:8:\the\ht0:\the\dp0}
\setbox0=\vbox{\hsize=1.490cm\parindent=0pt\linespread{1}\fontsize{10}{13.50}\selectfont\centering\noindent 播放\par}
\typeout{TCELL:p355_1:9:\the\ht0:\the\dp0}
\setbox0=\vbox{\hsize=4.940cm\parindent=0pt\linespread{1}\fontsize{10}{13.50}\selectfont\raggedright\noindent 将侵权作品的内容展示给不特定人或多数人，作品载体不转让出去\par}
\typeout{TCELL:p355_1:10:\the\ht0:\the\dp0}
\setbox0=\vbox{\hsize=3.540cm\parindent=0pt\linespread{1}\fontsize{10}{13.50}\selectfont\raggedright\noindent 线上播放，构成本罪（例3）\par}
\typeout{TCELL:p355_1:11:\the\ht0:\the\dp0}
\setbox0=\vbox{\hsize=3.590cm\parindent=0pt\linespread{1}\fontsize{10}{13.50}\selectfont\raggedright\noindent 线下播放，不构成本罪(例4)\par}
\typeout{TCELL:p355_1:12:\the\ht0:\the\dp0}
\setbox0=\vbox{\hsize=4.840cm\parindent=0pt\linespread{1}\fontsize{10}{13.50}\selectfont\bfseries\centering\noindent 本犯（甲）\par}
\typeout{TCELL:p374_1:0:\the\ht0:\the\dp0}
\setbox0=\vbox{\hsize=4.140cm\parindent=0pt\linespread{1}\fontsize{10}{13.50}\selectfont\bfseries\centering\noindent 帮助者（乙）\par}
\typeout{TCELL:p374_1:1:\the\ht0:\the\dp0}
\setbox0=\vbox{\hsize=6.240cm\parindent=0pt\linespread{1}\fontsize{10}{13.50}\selectfont\bfseries\centering\noindent 结论\par}
\typeout{TCELL:p374_1:2:\the\ht0:\the\dp0}
\setbox0=\vbox{\hsize=4.840cm\parindent=0pt\linespread{1}\fontsize{10}{13.50}\selectfont\raggedright\noindent 甲唆使乙帮助自己毁灭尸体\par}
\typeout{TCELL:p374_1:3:\the\ht0:\the\dp0}
\setbox0=\vbox{\hsize=4.140cm\parindent=0pt\linespread{1}\fontsize{10}{13.50}\selectfont\centering\noindent 乙毁灭尸体（实行）\par}
\typeout{TCELL:p374_1:4:\the\ht0:\the\dp0}
\setbox0=\vbox{\hsize=6.240cm\parindent=0pt\linespread{1}\fontsize{10}{13.50}\selectfont\raggedright\noindent 乙构成帮助毁灭证据罪；多数说认为，甲不构成该罪的教唆犯\par}
\typeout{TCELL:p374_1:5:\the\ht0:\the\dp0}
\setbox0=\vbox{\hsize=4.840cm\parindent=0pt\linespread{1}\fontsize{10}{13.50}\selectfont\raggedright\noindent 甲自己毁灭尸体，请乙提供刀具\par}
\typeout{TCELL:p374_1:6:\the\ht0:\the\dp0}
\setbox0=\vbox{\hsize=4.140cm\parindent=0pt\linespread{1}\fontsize{10}{13.50}\selectfont\centering\noindent 乙提供了刀具（帮助）\par}
\typeout{TCELL:p374_1:7:\the\ht0:\the\dp0}
\setbox0=\vbox{\hsize=6.240cm\parindent=0pt\linespread{1}\fontsize{10}{13.50}\selectfont\raggedright\noindent 乙构成帮助毁灭证据罪；甲不构成犯罪\par}
\typeout{TCELL:p374_1:8:\the\ht0:\the\dp0}
\setbox0=\vbox{\hsize=4.840cm\parindent=0pt\linespread{1}\fontsize{10}{13.50}\selectfont\raggedright\noindent 甲在乙的教唆下毁灭尸体\par}
\typeout{TCELL:p374_1:9:\the\ht0:\the\dp0}
\setbox0=\vbox{\hsize=4.140cm\parindent=0pt\linespread{1}\fontsize{10}{13.50}\selectfont\centering\noindent 乙教唆甲毁灭尸体（教唆）\par}
\typeout{TCELL:p374_1:10:\the\ht0:\the\dp0}
\setbox0=\vbox{\hsize=6.240cm\parindent=0pt\linespread{1}\fontsize{10}{13.50}\selectfont\raggedright\noindent 乙构成帮助毁灭证据罪；甲不构成犯罪\par}
\typeout{TCELL:p374_1:11:\the\ht0:\the\dp0}
\setbox0=\vbox{\hsize=1.840cm\parindent=0pt\linespread{1}\fontsize{10}{13.50}\selectfont\bfseries\centering\noindent 虚假诉讼罪\par}
\typeout{TCELL:p374_2:0:\the\ht0:\the\dp0}
\setbox0=\vbox{\hsize=3.640cm\parindent=0pt\linespread{1}\fontsize{10}{13.50}\selectfont\bfseries\centering\noindent 内容\par}
\typeout{TCELL:p374_2:1:\the\ht0:\the\dp0}
\setbox0=\vbox{\hsize=9.740cm\parindent=0pt\linespread{1}\fontsize{10}{13.50}\selectfont\bfseries\centering\noindent 具体表现\par}
\typeout{TCELL:p374_2:2:\the\ht0:\the\dp0}
\setbox0=\vbox{\hsize=1.840cm\parindent=0pt\linespread{1}\fontsize{10}{13.50}\selectfont\centering\noindent 行为方式（成立条件）\par}
\typeout{TCELL:p374_2:3:\the\ht0:\the\dp0}
\setbox0=\vbox{\hsize=3.640cm\parindent=0pt\linespread{1}\fontsize{10}{13.50}\selectfont\raggedright\noindent 以捏造的事实提起虚假的民事诉讼。\\{}1.执行程序属于诉讼的一部分；\\{}2.仲裁不属于诉讼\par}
\typeout{TCELL:p374_2:4:\the\ht0:\the\dp0}
\setbox0=\vbox{\hsize=9.740cm\parindent=0pt\linespread{1}\fontsize{10}{13.50}\selectfont\raggedright\noindent 1.乙欠甲30万元，甲将欠条修改为“欠80万元”，提起诉讼。多数说认为，“虚假”要求完全假，甲不构成虚假诉讼罪。少数说认为，“虚假”只要求部分假，甲构成虚假诉讼罪。\\{}2.隐瞒债务已经全部清偿的事实，向法院提起诉讼，要求还债。这种隐瞒事实，视为捏造事实，构成虚假诉讼罪\par}
\typeout{TCELL:p374_2:5:\the\ht0:\the\dp0}
\setbox0=\vbox{\hsize=1.840cm\parindent=0pt\linespread{1}\fontsize{10}{13.50}\selectfont\centering\noindent 行为后果（既遂条件）\par}
\typeout{TCELL:p374_2:6:\the\ht0:\the\dp0}
\setbox0=\vbox{\hsize=3.640cm\parindent=0pt\linespread{1}\fontsize{10}{13.50}\selectfont\raggedright\noindent 妨害司法秩序或者严重侵害他人合法权益\par}
\typeout{TCELL:p374_2:7:\the\ht0:\the\dp0}
\setbox0=\vbox{\hsize=9.740cm\parindent=0pt\linespread{1}\fontsize{10}{13.50}\selectfont\raggedright\noindent 1.致使法院基于捏造的事实\CJKunderwave{采取财产保全或者行为保全措施}的；\\{}2.\CJKunderwave{致使法院开庭审理}。\\{}满足任一条件，便构成既遂。（2023年试题）\\{}（少数说：法院受理便既遂）\par}
\typeout{TCELL:p374_2:8:\the\ht0:\the\dp0}
\setbox0=\vbox{\hsize=1.840cm\parindent=0pt\linespread{1}\fontsize{10}{13.50}\selectfont\centering\noindent 罪数\par}
\typeout{TCELL:p374_2:9:\the\ht0:\the\dp0}
\setbox0=\vbox{\hsize=13.640cm\parindent=0pt\linespread{1}\fontsize{10}{13.50}\selectfont\raggedright\noindent 1.甲以非法占有乙的财物为目的，提起虚假民事诉讼，法官受骗，判决甲胜诉，乙遭受财产损失。甲的一个行为同时触犯虚假诉讼罪和诈骗罪，想象竞合，择一重罪论处。\\{}2.甲与法官勾结，以非法占有乙的财物为目的，提起虚假民事诉讼，法官故意判决甲胜诉，乙遭受财产损失。甲与法官触犯三个罪的共同犯罪：针对司法秩序法益，触犯虚假诉讼罪、民事枉法裁判罪；针对被害人乙的财产，触犯盗窃罪（因为非法转移占有乙的财物）。想象竞合，择一重罪论处（2021年试题）\par}
\typeout{TCELL:p374_2:10:\the\ht0:\the\dp0}
\setbox0=\vbox{\hsize=4.690cm\parindent=0pt\linespread{1}\fontsize{10}{13.50}\selectfont\bfseries\centering\noindent 案件模型\par}
\typeout{TCELL:p379_1:0:\the\ht0:\the\dp0}
\setbox0=\vbox{\hsize=4.690cm\parindent=0pt\linespread{1}\fontsize{10}{13.50}\selectfont\bfseries\centering\noindent 正常销赃，未欺骗买家\par}
\typeout{TCELL:p379_1:1:\the\ht0:\the\dp0}
\setbox0=\vbox{\hsize=5.840cm\parindent=0pt\linespread{1}\fontsize{10}{13.50}\selectfont\bfseries\centering\noindent 欺骗买家，谎称是自己的财物\par}
\typeout{TCELL:p379_1:2:\the\ht0:\the\dp0}
\setbox0=\vbox{\hsize=4.690cm\parindent=0pt\linespread{1}\fontsize{10}{13.50}\selectfont\centering\noindent 甲偷了丙的财物，甲销赃\par}
\typeout{TCELL:p379_1:3:\the\ht0:\the\dp0}
\setbox0=\vbox{\hsize=4.690cm\parindent=0pt\linespread{1}\fontsize{10}{13.50}\selectfont\raggedright\noindent 不可罚的事后行为\par}
\typeout{TCELL:p379_1:4:\the\ht0:\the\dp0}
\setbox0=\vbox{\hsize=5.840cm\parindent=0pt\linespread{1}\fontsize{10}{13.50}\selectfont\raggedright\noindent 定诈骗罪，与盗窃罪并罚\par}
\typeout{TCELL:p379_1:5:\the\ht0:\the\dp0}
\setbox0=\vbox{\hsize=4.690cm\parindent=0pt\linespread{1}\fontsize{10}{13.50}\selectfont\centering\noindent 甲偷了丙的财物，乙帮甲销赃\par}
\typeout{TCELL:p379_1:6:\the\ht0:\the\dp0}
\setbox0=\vbox{\hsize=4.690cm\parindent=0pt\linespread{1}\fontsize{10}{13.50}\selectfont\raggedright\noindent 定掩饰、隐瞒犯罪所得罪\par}
\typeout{TCELL:p379_1:7:\the\ht0:\the\dp0}
\setbox0=\vbox{\hsize=5.840cm\parindent=0pt\linespread{1}\fontsize{10}{13.50}\selectfont\raggedright\noindent 同时触犯掩饰、隐瞒犯罪所得罪和诈骗罪，想象竞合，择一重罪论处\par}
\typeout{TCELL:p379_1:8:\the\ht0:\the\dp0}
\setbox0=\vbox{\hsize=1.540cm\parindent=0pt\linespread{1}\fontsize{10}{13.50}\selectfont\bfseries\centering\noindent 行为类型\par}
\typeout{TCELL:p386_2:0:\the\ht0:\the\dp0}
\setbox0=\vbox{\hsize=8.040cm\parindent=0pt\linespread{1}\fontsize{10}{13.50}\selectfont\bfseries\centering\noindent 具体模型\par}
\typeout{TCELL:p386_2:1:\the\ht0:\the\dp0}
\setbox0=\vbox{\hsize=5.640cm\parindent=0pt\linespread{1}\fontsize{10}{13.50}\selectfont\bfseries\centering\noindent 结论\par}
\typeout{TCELL:p386_2:2:\the\ht0:\the\dp0}
\setbox0=\vbox{\hsize=1.540cm\parindent=0pt\linespread{1}\fontsize{10}{13.50}\selectfont\centering\noindent 居间介绍\par}
\typeout{TCELL:p386_2:3:\the\ht0:\the\dp0}
\setbox0=\vbox{\hsize=8.040cm\parindent=0pt\linespread{1}\fontsize{10}{13.50}\selectfont\raggedright\noindent 贩毒者乙想找货源，甲受乙委托，为乙介绍贩毒者\par}
\typeout{TCELL:p386_2:4:\the\ht0:\the\dp0}
\setbox0=\vbox{\hsize=8.040cm\parindent=0pt\linespread{1}\fontsize{10}{13.50}\selectfont\raggedright\noindent 贩毒者乙想找买家，甲受乙委托，为乙介绍购毒者\par}
\typeout{TCELL:p386_2:5:\the\ht0:\the\dp0}
\setbox0=\vbox{\hsize=5.640cm\parindent=0pt\linespread{1}\fontsize{10}{13.50}\selectfont\raggedright\noindent 甲构成乙的贩卖毒品罪的帮助犯\par}
\typeout{TCELL:p386_2:6:\the\ht0:\the\dp0}
\setbox0=\vbox{\hsize=8.040cm\parindent=0pt\linespread{1}\fontsize{10}{13.50}\selectfont\raggedright\noindent 吸毒者乙想买毒品，甲受乙委托，为乙介绍贩毒者，乙买到数量较多毒品而持有\par}
\typeout{TCELL:p386_2:7:\the\ht0:\the\dp0}
\setbox0=\vbox{\hsize=5.640cm\parindent=0pt\linespread{1}\fontsize{10}{13.50}\selectfont\raggedright\noindent 甲构成乙的非法持有毒品罪的帮助犯\par}
\typeout{TCELL:p386_2:8:\the\ht0:\the\dp0}
\setbox0=\vbox{\hsize=8.040cm\parindent=0pt\linespread{1}\fontsize{10}{13.50}\selectfont\raggedright\noindent 甲在贩毒者与吸毒者之间居间介绍，促成交易\par}
\typeout{TCELL:p386_2:9:\the\ht0:\the\dp0}
\setbox0=\vbox{\hsize=5.640cm\parindent=0pt\linespread{1}\fontsize{10}{13.50}\selectfont\raggedright\noindent 甲一方面构成贩毒者的贩卖毒品罪的帮助犯，另一方面构成吸毒者的非法持有毒品罪的帮助犯，想象竞合\par}
\typeout{TCELL:p386_2:10:\the\ht0:\the\dp0}
\setbox0=\vbox{\hsize=1.540cm\parindent=0pt\linespread{1}\fontsize{10}{13.50}\selectfont\bfseries\centering\noindent 行为类型\par}
\typeout{TCELL:p387_1:0:\the\ht0:\the\dp0}
\setbox0=\vbox{\hsize=8.040cm\parindent=0pt\linespread{1}\fontsize{10}{13.50}\selectfont\bfseries\centering\noindent 具体模型\par}
\typeout{TCELL:p387_1:1:\the\ht0:\the\dp0}
\setbox0=\vbox{\hsize=5.640cm\parindent=0pt\linespread{1}\fontsize{10}{13.50}\selectfont\bfseries\centering\noindent 结论\par}
\typeout{TCELL:p387_1:2:\the\ht0:\the\dp0}
\setbox0=\vbox{\hsize=1.540cm\parindent=0pt\linespread{1}\fontsize{10}{13.50}\selectfont\centering\noindent 居间倒卖\par}
\typeout{TCELL:p387_1:3:\the\ht0:\the\dp0}
\setbox0=\vbox{\hsize=8.040cm\parindent=0pt\linespread{1}\fontsize{10}{13.50}\selectfont\raggedright\noindent 贩毒者将毒品卖给甲，甲将毒品倒卖给吸毒者\par}
\typeout{TCELL:p387_1:4:\the\ht0:\the\dp0}
\setbox0=\vbox{\hsize=5.640cm\parindent=0pt\linespread{1}\fontsize{10}{13.50}\selectfont\raggedright\noindent 甲构成贩卖毒品罪的实行犯\par}
\typeout{TCELL:p387_1:5:\the\ht0:\the\dp0}
\setbox0=\vbox{\hsize=1.540cm\parindent=0pt\linespread{1}\fontsize{10}{13.50}\selectfont\centering\noindent 帮助交付\par}
\typeout{TCELL:p387_1:6:\the\ht0:\the\dp0}
\setbox0=\vbox{\hsize=8.040cm\parindent=0pt\linespread{1}\fontsize{10}{13.50}\selectfont\raggedright\noindent 贩毒者将毒品卖给吸毒者，让甲交付给吸毒者\par}
\typeout{TCELL:p387_1:7:\the\ht0:\the\dp0}
\setbox0=\vbox{\hsize=5.640cm\parindent=0pt\linespread{1}\fontsize{10}{13.50}\selectfont\raggedright\noindent 甲构成贩卖毒品罪的帮助犯\par}
\typeout{TCELL:p387_1:8:\the\ht0:\the\dp0}
\setbox0=\vbox{\hsize=3.240cm\parindent=0pt\linespread{1}\fontsize{10}{13.50}\selectfont\bfseries\centering\noindent 角色与行为\par}
\typeout{TCELL:p387_2:0:\the\ht0:\the\dp0}
\setbox0=\vbox{\hsize=6.690cm\parindent=0pt\linespread{1}\fontsize{10}{13.50}\selectfont\bfseries\centering\noindent 购买、存储毒品过程中，持有毒品数量较大\par}
\typeout{TCELL:p387_2:1:\the\ht0:\the\dp0}
\setbox0=\vbox{\hsize=5.290cm\parindent=0pt\linespread{1}\fontsize{10}{13.50}\selectfont\bfseries\centering\noindent 运输毒品过程中，\CJKunderwave{运输毒品数量}\CJKunderwave{较大}\par}
\typeout{TCELL:p387_2:2:\the\ht0:\the\dp0}
\setbox0=\vbox{\hsize=3.240cm\parindent=0pt\linespread{1}\fontsize{10}{13.50}\selectfont\centering\noindent 吸毒者，代购者（为吸毒者代购）\par}
\typeout{TCELL:p387_2:3:\the\ht0:\the\dp0}
\setbox0=\vbox{\hsize=6.690cm\parindent=0pt\linespread{1}\fontsize{10}{13.50}\selectfont\raggedright\noindent 非法持有毒品罪\par}
\typeout{TCELL:p387_2:4:\the\ht0:\the\dp0}
\setbox0=\vbox{\hsize=5.290cm\parindent=0pt\linespread{1}\fontsize{10}{13.50}\selectfont\raggedright\noindent 运输毒品罪\par}
\typeout{TCELL:p387_2:5:\the\ht0:\the\dp0}
\setbox0=\vbox{\hsize=2.190cm\parindent=0pt\linespread{1}\fontsize{10}{13.50}\selectfont\bfseries\centering\noindent 组织卖淫罪\par}
\typeout{TCELL:p393_1:0:\the\ht0:\the\dp0}
\setbox0=\vbox{\hsize=4.740cm\parindent=0pt\linespread{1}\fontsize{10}{13.50}\selectfont\bfseries\centering\noindent 引诱、容留、介绍卖淫罪\par}
\typeout{TCELL:p393_1:1:\the\ht0:\the\dp0}
\setbox0=\vbox{\hsize=8.290cm\parindent=0pt\linespread{1}\fontsize{10}{13.50}\selectfont\bfseries\centering\noindent 处理结论\par}
\typeout{TCELL:p393_1:2:\the\ht0:\the\dp0}
\setbox0=\vbox{\hsize=2.190cm\parindent=0pt\linespread{1}\fontsize{10}{13.50}\selectfont\raggedright\noindent \par}
\typeout{TCELL:p393_1:3:\the\ht0:\the\dp0}
\setbox0=\vbox{\hsize=4.740cm\parindent=0pt\linespread{1}\fontsize{10}{13.50}\selectfont\raggedright\noindent 引诱甲、容留乙、介绍丙\par}
\typeout{TCELL:p393_1:4:\the\ht0:\the\dp0}
\setbox0=\vbox{\hsize=8.290cm\parindent=0pt\linespread{1}\fontsize{10}{13.50}\selectfont\raggedright\noindent 定一个引诱、容留、介绍卖淫罪（选择性罪名，对象虽不同一，但不并罚)\par}
\typeout{TCELL:p393_1:5:\the\ht0:\the\dp0}
\setbox0=\vbox{\hsize=2.190cm\parindent=0pt\linespread{1}\fontsize{10}{13.50}\selectfont\centering\noindent 组织甲\par}
\typeout{TCELL:p393_1:6:\the\ht0:\the\dp0}
\setbox0=\vbox{\hsize=4.740cm\parindent=0pt\linespread{1}\fontsize{10}{13.50}\selectfont\raggedright\noindent 并且，引诱甲、容留甲、介绍甲\par}
\typeout{TCELL:p393_1:7:\the\ht0:\the\dp0}
\setbox0=\vbox{\hsize=8.290cm\parindent=0pt\linespread{1}\fontsize{10}{13.50}\selectfont\raggedright\noindent 定一个组织卖淫罪（两个罪名，对象具有同一性，不需并罚，前者吸收后者）\par}
\typeout{TCELL:p393_1:8:\the\ht0:\the\dp0}
\setbox0=\vbox{\hsize=2.190cm\parindent=0pt\linespread{1}\fontsize{10}{13.50}\selectfont\centering\noindent 组织甲\par}
\typeout{TCELL:p393_1:9:\the\ht0:\the\dp0}
\setbox0=\vbox{\hsize=4.740cm\parindent=0pt\linespread{1}\fontsize{10}{13.50}\selectfont\raggedright\noindent 并且，引诱乙、容留乙、介绍乙\par}
\typeout{TCELL:p393_1:10:\the\ht0:\the\dp0}
\setbox0=\vbox{\hsize=8.290cm\parindent=0pt\linespread{1}\fontsize{10}{13.50}\selectfont\raggedright\noindent 定组织卖淫罪和引诱、容留、介绍卖淫罪，并罚（两个罪名，对象不具有同一性，需并罚）\par}
\typeout{TCELL:p393_1:11:\the\ht0:\the\dp0}
\setbox0=\vbox{\hsize=2.190cm\parindent=0pt\linespread{1}\fontsize{10}{13.50}\selectfont\raggedright\noindent \par}
\typeout{TCELL:p393_1:12:\the\ht0:\the\dp0}
\setbox0=\vbox{\hsize=4.740cm\parindent=0pt\linespread{1}\fontsize{10}{13.50}\selectfont\raggedright\noindent 引诱、容留、介绍幼女卖淫\par}
\typeout{TCELL:p393_1:13:\the\ht0:\the\dp0}
\setbox0=\vbox{\hsize=8.290cm\parindent=0pt\linespread{1}\fontsize{10}{13.50}\selectfont\raggedright\noindent 引诱幼女卖淫罪与容留、介绍卖淫罪，并罚（2011年第56题）\par}
\typeout{TCELL:p393_1:14:\the\ht0:\the\dp0}
\setbox0=\vbox{\hsize=7.290cm\parindent=0pt\linespread{1}\fontsize{10}{13.50}\selectfont\bfseries\centering\noindent 罪名\par}
\typeout{TCELL:p395_1:0:\the\ht0:\the\dp0}
\setbox0=\vbox{\hsize=8.190cm\parindent=0pt\linespread{1}\fontsize{10}{13.50}\selectfont\bfseries\centering\noindent 要求的目的\par}
\typeout{TCELL:p395_1:1:\the\ht0:\the\dp0}
\setbox0=\vbox{\hsize=7.290cm\parindent=0pt\linespread{1}\fontsize{10}{13.50}\selectfont\centering\noindent 传播淫秽物品罪\par}
\typeout{TCELL:p395_1:2:\the\ht0:\the\dp0}
\setbox0=\vbox{\hsize=8.190cm\parindent=0pt\linespread{1}\fontsize{10}{13.50}\selectfont\raggedright\noindent 传播目的\par}
\typeout{TCELL:p395_1:3:\the\ht0:\the\dp0}
\setbox0=\vbox{\hsize=7.290cm\parindent=0pt\linespread{1}\fontsize{10}{13.50}\selectfont\centering\noindent 传播淫秽物品牟利罪\par}
\typeout{TCELL:p395_1:4:\the\ht0:\the\dp0}
\setbox0=\vbox{\hsize=8.190cm\parindent=0pt\linespread{1}\fontsize{10}{13.50}\selectfont\raggedright\noindent 传播目的，并且有牟利目的\par}
\typeout{TCELL:p395_1:5:\the\ht0:\the\dp0}
\setbox0=\vbox{\hsize=7.290cm\parindent=0pt\linespread{1}\fontsize{10}{13.50}\selectfont\centering\noindent 走私淫秽物品罪（第152条）\par}
\typeout{TCELL:p395_1:6:\the\ht0:\the\dp0}
\setbox0=\vbox{\hsize=8.190cm\parindent=0pt\linespread{1}\fontsize{10}{13.50}\selectfont\raggedright\noindent 传播目的，或者有牟利目的（2015年第61题）\par}
\typeout{TCELL:p395_1:7:\the\ht0:\the\dp0}
\setbox0=\vbox{\hsize=3.090cm\parindent=0pt\linespread{1}\fontsize{10}{13.50}\selectfont\bfseries\centering\noindent 盗窃对象\par}
\typeout{TCELL:p395_2:0:\the\ht0:\the\dp0}
\setbox0=\vbox{\hsize=2.790cm\parindent=0pt\linespread{1}\fontsize{10}{13.50}\selectfont\bfseries\centering\noindent 定的罪名\par}
\typeout{TCELL:p395_2:1:\the\ht0:\the\dp0}
\setbox0=\vbox{\hsize=4.840cm\parindent=0pt\linespread{1}\fontsize{10}{13.50}\selectfont\bfseries\centering\noindent 又销售，定的罪名\par}
\typeout{TCELL:p395_2:2:\the\ht0:\the\dp0}
\setbox0=\vbox{\hsize=4.240cm\parindent=0pt\linespread{1}\fontsize{10}{13.50}\selectfont\bfseries\centering\noindent 罪数处理\par}
\typeout{TCELL:p395_2:3:\the\ht0:\the\dp0}
\setbox0=\vbox{\hsize=3.090cm\parindent=0pt\linespread{1}\fontsize{10}{13.50}\selectfont\centering\noindent 假币\par}
\typeout{TCELL:p395_2:4:\the\ht0:\the\dp0}
\setbox0=\vbox{\hsize=2.790cm\parindent=0pt\linespread{1}\fontsize{10}{13.50}\selectfont\centering\noindent 盗窃罪\par}
\typeout{TCELL:p395_2:5:\the\ht0:\the\dp0}
\setbox0=\vbox{\hsize=4.840cm\parindent=0pt\linespread{1}\fontsize{10}{13.50}\selectfont\centering\noindent 出售假币罪\par}
\typeout{TCELL:p395_2:6:\the\ht0:\the\dp0}
\setbox0=\vbox{\hsize=4.240cm\parindent=0pt\linespread{1}\fontsize{10}{13.50}\selectfont\centering\noindent 前后两个行为，构成两个罪，数罪并罚\par}
\typeout{TCELL:p395_2:7:\the\ht0:\the\dp0}
\setbox0=\vbox{\hsize=3.090cm\parindent=0pt\linespread{1}\fontsize{10}{13.50}\selectfont\centering\noindent 淫秽物品\par}
\typeout{TCELL:p395_2:8:\the\ht0:\the\dp0}
\setbox0=\vbox{\hsize=4.840cm\parindent=0pt\linespread{1}\fontsize{10}{13.50}\selectfont\centering\noindent 贩卖淫秽物品牟利罪\par}
\typeout{TCELL:p395_2:9:\the\ht0:\the\dp0}
\setbox0=\vbox{\hsize=3.090cm\parindent=0pt\linespread{1}\fontsize{10}{13.50}\selectfont\centering\noindent 文物\par}
\typeout{TCELL:p395_2:10:\the\ht0:\the\dp0}
\setbox0=\vbox{\hsize=4.840cm\parindent=0pt\linespread{1}\fontsize{10}{13.50}\selectfont\centering\noindent 倒卖文物罪\par}
\typeout{TCELL:p395_2:11:\the\ht0:\the\dp0}
\setbox0=\vbox{\hsize=3.090cm\parindent=0pt\linespread{1}\fontsize{10}{13.50}\selectfont\centering\noindent 毒品\par}
\typeout{TCELL:p395_2:12:\the\ht0:\the\dp0}
\setbox0=\vbox{\hsize=4.840cm\parindent=0pt\linespread{1}\fontsize{10}{13.50}\selectfont\centering\noindent 贩卖毒品罪\par}
\typeout{TCELL:p395_2:13:\the\ht0:\the\dp0}
\setbox0=\vbox{\hsize=3.090cm\parindent=0pt\linespread{1}\fontsize{10}{13.50}\selectfont\centering\noindent 枪支\par}
\typeout{TCELL:p395_2:14:\the\ht0:\the\dp0}
\setbox0=\vbox{\hsize=2.790cm\parindent=0pt\linespread{1}\fontsize{10}{13.50}\selectfont\centering\noindent 盗窃枪支罪\par}
\typeout{TCELL:p395_2:15:\the\ht0:\the\dp0}
\setbox0=\vbox{\hsize=4.840cm\parindent=0pt\linespread{1}\fontsize{10}{13.50}\selectfont\centering\noindent 非法买卖枪支罪\par}
\typeout{TCELL:p395_2:16:\the\ht0:\the\dp0}
\setbox0=\vbox{\hsize=4.240cm\parindent=0pt\linespread{1}\fontsize{10}{13.50}\selectfont\centering\noindent 只定盗窃枪支罪\par}
\typeout{TCELL:p395_2:17:\the\ht0:\the\dp0}
\setbox0=\vbox{\hsize=2.390cm\parindent=0pt\linespread{1}\fontsize{10}{13.50}\selectfont\bfseries\centering\noindent 走私对象\par}
\typeout{TCELL:p395_3:0:\the\ht0:\the\dp0}
\setbox0=\vbox{\hsize=4.140cm\parindent=0pt\linespread{1}\fontsize{10}{13.50}\selectfont\bfseries\centering\noindent 定的罪名\par}
\typeout{TCELL:p395_3:1:\the\ht0:\the\dp0}
\setbox0=\vbox{\hsize=3.940cm\parindent=0pt\linespread{1}\fontsize{10}{13.50}\selectfont\bfseries\centering\noindent 又销售，定的罪名\par}
\typeout{TCELL:p395_3:2:\the\ht0:\the\dp0}
\setbox0=\vbox{\hsize=4.490cm\parindent=0pt\linespread{1}\fontsize{10}{13.50}\selectfont\bfseries\centering\noindent 罪数处理\par}
\typeout{TCELL:p395_3:3:\the\ht0:\the\dp0}
\setbox0=\vbox{\hsize=2.390cm\parindent=0pt\linespread{1}\fontsize{10}{13.50}\selectfont\centering\noindent 假币\par}
\typeout{TCELL:p395_3:4:\the\ht0:\the\dp0}
\setbox0=\vbox{\hsize=4.140cm\parindent=0pt\linespread{1}\fontsize{10}{13.50}\selectfont\centering\noindent 走私假币罪\par}
\typeout{TCELL:p395_3:5:\the\ht0:\the\dp0}
\setbox0=\vbox{\hsize=3.940cm\parindent=0pt\linespread{1}\fontsize{10}{13.50}\selectfont\centering\noindent 出售假币罪\par}
\typeout{TCELL:p395_3:6:\the\ht0:\the\dp0}
\setbox0=\vbox{\hsize=4.490cm\parindent=0pt\linespread{1}\fontsize{10}{13.50}\selectfont\centering\noindent 只定前面的走私犯罪\par}
\typeout{TCELL:p395_3:7:\the\ht0:\the\dp0}
\setbox0=\vbox{\hsize=2.390cm\parindent=0pt\linespread{1}\fontsize{10}{13.50}\selectfont\centering\noindent 淫秽物品\par}
\typeout{TCELL:p395_3:8:\the\ht0:\the\dp0}
\setbox0=\vbox{\hsize=4.140cm\parindent=0pt\linespread{1}\fontsize{10}{13.50}\selectfont\centering\noindent 走私淫秽物品罪\par}
\typeout{TCELL:p395_3:9:\the\ht0:\the\dp0}
\setbox0=\vbox{\hsize=3.940cm\parindent=0pt\linespread{1}\fontsize{10}{13.50}\selectfont\centering\noindent 贩卖淫秽物品牟利罪\par}
\typeout{TCELL:p395_3:10:\the\ht0:\the\dp0}
\setbox0=\vbox{\hsize=2.390cm\parindent=0pt\linespread{1}\fontsize{10}{13.50}\selectfont\centering\noindent 文物\par}
\typeout{TCELL:p395_3:11:\the\ht0:\the\dp0}
\setbox0=\vbox{\hsize=4.140cm\parindent=0pt\linespread{1}\fontsize{10}{13.50}\selectfont\raggedright\noindent 入境，不定走私文物罪；\\{}出境，定走私文物罪\par}
\typeout{TCELL:p395_3:12:\the\ht0:\the\dp0}
\setbox0=\vbox{\hsize=3.940cm\parindent=0pt\linespread{1}\fontsize{10}{13.50}\selectfont\raggedright\noindent 倒卖文物罪，非法向外国人出售珍贵文物罪\par}
\typeout{TCELL:p395_3:13:\the\ht0:\the\dp0}
\setbox0=\vbox{\hsize=2.390cm\parindent=0pt\linespread{1}\fontsize{10}{13.50}\selectfont\centering\noindent 毒品\par}
\typeout{TCELL:p395_3:14:\the\ht0:\the\dp0}
\setbox0=\vbox{\hsize=4.140cm\parindent=0pt\linespread{1}\fontsize{10}{13.50}\selectfont\centering\noindent 走私毒品罪\par}
\typeout{TCELL:p395_3:15:\the\ht0:\the\dp0}
\setbox0=\vbox{\hsize=3.940cm\parindent=0pt\linespread{1}\fontsize{10}{13.50}\selectfont\centering\noindent 贩卖毒品罪\par}
\typeout{TCELL:p395_3:16:\the\ht0:\the\dp0}
\setbox0=\vbox{\hsize=4.490cm\parindent=0pt\linespread{1}\fontsize{10}{13.50}\selectfont\raggedright\noindent 定走私、贩卖毒品罪，不并罚\par}
\typeout{TCELL:p395_3:17:\the\ht0:\the\dp0}
\setbox0=\vbox{\hsize=2.390cm\parindent=0pt\linespread{1}\fontsize{10}{13.50}\selectfont\centering\noindent 枪支\par}
\typeout{TCELL:p395_3:18:\the\ht0:\the\dp0}
\setbox0=\vbox{\hsize=4.140cm\parindent=0pt\linespread{1}\fontsize{10}{13.50}\selectfont\centering\noindent 走私武器罪\par}
\typeout{TCELL:p395_3:19:\the\ht0:\the\dp0}
\setbox0=\vbox{\hsize=3.940cm\parindent=0pt\linespread{1}\fontsize{10}{13.50}\selectfont\centering\noindent 非法买卖枪支罪\par}
\typeout{TCELL:p395_3:20:\the\ht0:\the\dp0}
\setbox0=\vbox{\hsize=4.490cm\parindent=0pt\linespread{1}\fontsize{10}{13.50}\selectfont\raggedright\noindent 两罪并罚（因为后罪是危害公共安全犯罪，危害性很大，前罪吸收不了后罪）\par}
\typeout{TCELL:p395_3:21:\the\ht0:\the\dp0}
\setbox0=\vbox{\hsize=3.940cm\parindent=0pt\linespread{1}\fontsize{10}{13.50}\selectfont\bfseries\centering\noindent 非国家工作人员（甲）\par}
\typeout{TCELL:p400_1:0:\the\ht0:\the\dp0}
\setbox0=\vbox{\hsize=3.290cm\parindent=0pt\linespread{1}\fontsize{10}{13.50}\selectfont\bfseries\centering\noindent 国家工作人员（乙）\par}
\typeout{TCELL:p400_1:1:\the\ht0:\the\dp0}
\setbox0=\vbox{\hsize=7.990cm\parindent=0pt\linespread{1}\fontsize{10}{13.50}\selectfont\bfseries\centering\noindent 结论\par}
\typeout{TCELL:p400_1:2:\the\ht0:\the\dp0}
\setbox0=\vbox{\hsize=3.940cm\parindent=0pt\linespread{1}\fontsize{10}{13.50}\selectfont\centering\noindent 帮助行为\par}
\typeout{TCELL:p400_1:3:\the\ht0:\the\dp0}
\setbox0=\vbox{\hsize=3.290cm\parindent=0pt\linespread{1}\fontsize{10}{13.50}\selectfont\raggedright\noindent 实行行为\par}
\typeout{TCELL:p400_1:4:\the\ht0:\the\dp0}
\setbox0=\vbox{\hsize=7.990cm\parindent=0pt\linespread{1}\fontsize{10}{13.50}\selectfont\raggedright\noindent 乙构成贪污罪（实行犯），甲构成贪污罪（帮助犯）\\{}（例1）\par}
\typeout{TCELL:p400_1:5:\the\ht0:\the\dp0}
\setbox0=\vbox{\hsize=3.940cm\parindent=0pt\linespread{1}\fontsize{10}{13.50}\selectfont\bfseries\centering\noindent 非国家工作人员（甲）\par}
\typeout{TCELL:p401_1:0:\the\ht0:\the\dp0}
\setbox0=\vbox{\hsize=3.290cm\parindent=0pt\linespread{1}\fontsize{10}{13.50}\selectfont\bfseries\centering\noindent 国家工作人员（乙）\par}
\typeout{TCELL:p401_1:1:\the\ht0:\the\dp0}
\setbox0=\vbox{\hsize=7.990cm\parindent=0pt\linespread{1}\fontsize{10}{13.50}\selectfont\bfseries\centering\noindent 结论\par}
\typeout{TCELL:p401_1:2:\the\ht0:\the\dp0}
\setbox0=\vbox{\hsize=3.940cm\parindent=0pt\linespread{1}\fontsize{10}{13.50}\selectfont\centering\noindent 实行行为\par}
\typeout{TCELL:p401_1:3:\the\ht0:\the\dp0}
\setbox0=\vbox{\hsize=3.290cm\parindent=0pt\linespread{1}\fontsize{10}{13.50}\selectfont\raggedright\noindent 实行行为\par}
\typeout{TCELL:p401_1:4:\the\ht0:\the\dp0}
\setbox0=\vbox{\hsize=7.990cm\parindent=0pt\linespread{1}\fontsize{10}{13.50}\selectfont\raggedright\noindent 甲乙既构成财产犯罪的共同犯罪，也构成贪污罪的共同犯罪，想象竞合，择一重罪论处（例2）\par}
\typeout{TCELL:p401_1:5:\the\ht0:\the\dp0}
\setbox0=\vbox{\hsize=3.940cm\parindent=0pt\linespread{1}\fontsize{10}{13.50}\selectfont\centering\noindent 实行行为\par}
\typeout{TCELL:p401_1:6:\the\ht0:\the\dp0}
\setbox0=\vbox{\hsize=3.290cm\parindent=0pt\linespread{1}\fontsize{10}{13.50}\selectfont\raggedright\noindent 帮助行为\par}
\typeout{TCELL:p401_1:7:\the\ht0:\the\dp0}
\setbox0=\vbox{\hsize=7.990cm\parindent=0pt\linespread{1}\fontsize{10}{13.50}\selectfont\raggedright\noindent 甲构成财产犯罪（实行犯），乙构成财产犯罪（帮助犯）。乙不构成贪污罪（例3）\par}
\typeout{TCELL:p401_1:8:\the\ht0:\the\dp0}
\setbox0=\vbox{\hsize=1.640cm\parindent=0pt\linespread{1}\fontsize{10}{13.50}\selectfont\bfseries\centering\noindent 罪名\par}
\typeout{TCELL:p406_1:0:\the\ht0:\the\dp0}
\setbox0=\vbox{\hsize=2.940cm\parindent=0pt\linespread{1}\fontsize{10}{13.50}\selectfont\bfseries\centering\noindent 挪用资金罪\par}
\typeout{TCELL:p406_1:1:\the\ht0:\the\dp0}
\setbox0=\vbox{\hsize=2.540cm\parindent=0pt\linespread{1}\fontsize{10}{13.50}\selectfont\bfseries\centering\noindent 挪用公款罪\par}
\typeout{TCELL:p406_1:2:\the\ht0:\the\dp0}
\setbox0=\vbox{\hsize=7.840cm\parindent=0pt\linespread{1}\fontsize{10}{13.50}\selectfont\bfseries\centering\noindent 挪用特定款物罪\par}
\typeout{TCELL:p406_1:3:\the\ht0:\the\dp0}
\setbox0=\vbox{\hsize=1.640cm\parindent=0pt\linespread{1}\fontsize{10}{13.50}\selectfont\centering\noindent 行为主体\par}
\typeout{TCELL:p406_1:4:\the\ht0:\the\dp0}
\setbox0=\vbox{\hsize=2.940cm\parindent=0pt\linespread{1}\fontsize{10}{13.50}\selectfont\centering\noindent 非国家工作人员\par}
\typeout{TCELL:p406_1:5:\the\ht0:\the\dp0}
\setbox0=\vbox{\hsize=2.540cm\parindent=0pt\linespread{1}\fontsize{10}{13.50}\selectfont\centering\noindent 国家工作人员\par}
\typeout{TCELL:p406_1:6:\the\ht0:\the\dp0}
\setbox0=\vbox{\hsize=7.840cm\parindent=0pt\linespread{1}\fontsize{10}{13.50}\selectfont\raggedright\noindent 一般主体，包括国家工作人员和非国家工作人员①\par}
\typeout{TCELL:p406_1:7:\the\ht0:\the\dp0}
\setbox0=\vbox{\hsize=1.640cm\parindent=0pt\linespread{1}\fontsize{10}{13.50}\selectfont\centering\noindent 行为对象\par}
\typeout{TCELL:p406_1:8:\the\ht0:\the\dp0}
\setbox0=\vbox{\hsize=2.940cm\parindent=0pt\linespread{1}\fontsize{10}{13.50}\selectfont\centering\noindent 非国有单位的资金\par}
\typeout{TCELL:p406_1:9:\the\ht0:\the\dp0}
\setbox0=\vbox{\hsize=2.540cm\parindent=0pt\linespread{1}\fontsize{10}{13.50}\selectfont\centering\noindent 公家的资金\par}
\typeout{TCELL:p406_1:10:\the\ht0:\the\dp0}
\setbox0=\vbox{\hsize=7.840cm\parindent=0pt\linespread{1}\fontsize{10}{13.50}\selectfont\raggedright\noindent 特定款物，指专用于抢险、救灾、救济等用途的款（资金）和物（物品）\par}
\typeout{TCELL:p406_1:11:\the\ht0:\the\dp0}
\setbox0=\vbox{\hsize=1.640cm\parindent=0pt\linespread{1}\fontsize{10}{13.50}\selectfont\centering\noindent 行为方式\par}
\typeout{TCELL:p406_1:12:\the\ht0:\the\dp0}
\setbox0=\vbox{\hsize=2.940cm\parindent=0pt\linespread{1}\fontsize{10}{13.50}\selectfont\raggedright\noindent 挪作私用，有归还意图\par}
\typeout{TCELL:p406_1:13:\the\ht0:\the\dp0}
\setbox0=\vbox{\hsize=2.540cm\parindent=0pt\linespread{1}\fontsize{10}{13.50}\selectfont\raggedright\noindent 挪作私用，有归还意图\par}
\typeout{TCELL:p406_1:14:\the\ht0:\the\dp0}
\setbox0=\vbox{\hsize=7.840cm\parindent=0pt\linespread{1}\fontsize{10}{13.50}\selectfont\raggedright\noindent 将A公用用途变成B公用用途，有无归还意图，在所不问\par}
\typeout{TCELL:p406_1:15:\the\ht0:\the\dp0}
\setbox0=\vbox{\hsize=1.640cm\parindent=0pt\linespread{1}\fontsize{10}{13.50}\selectfont\centering\noindent 要点\par}
\typeout{TCELL:p406_1:16:\the\ht0:\the\dp0}
\setbox0=\vbox{\hsize=13.840cm\parindent=0pt\linespread{1}\fontsize{10}{13.50}\selectfont\raggedright\noindent 1.将特定款物，挪作私用，定挪用公款罪。其中的“特定物”定挪用公款罪，属于特殊规定，因为“物”本不属于“款”。\\{}2.失业保险基金、下岗职工基本生活保障资金，属于救济类特定款物。根据行为方式不同，可定挪用特定款物罪或挪用公款罪。\\{}3.退休职工养老保险金，不属于特定款物。根据行为主体身份不同，可定挪用资金罪或挪用公款罪\par}
\typeout{TCELL:p406_1:17:\the\ht0:\the\dp0}
\setbox0=\vbox{\hsize=2.840cm\parindent=0pt\linespread{1}\fontsize{10}{13.50}\selectfont\bfseries\centering\noindent 误解的类型\par}
\typeout{TCELL:p411_1:0:\the\ht0:\the\dp0}
\setbox0=\vbox{\hsize=3.940cm\parindent=0pt\linespread{1}\fontsize{10}{13.50}\selectfont\bfseries\centering\noindent 具体情形\par}
\typeout{TCELL:p411_1:1:\the\ht0:\the\dp0}
\setbox0=\vbox{\hsize=8.440cm\parindent=0pt\linespread{1}\fontsize{10}{13.50}\selectfont\bfseries\centering\noindent 范例\par}
\typeout{TCELL:p411_1:2:\the\ht0:\the\dp0}
\setbox0=\vbox{\hsize=2.840cm\parindent=0pt\linespread{1}\fontsize{10}{13.50}\selectfont\raggedright\noindent 误解一：只要客观上具备数额，就定罪量刑\par}
\typeout{TCELL:p411_1:3:\the\ht0:\the\dp0}
\setbox0=\vbox{\hsize=3.940cm\parindent=0pt\linespread{1}\fontsize{10}{13.50}\selectfont\raggedright\noindent （1）只要客观上达到“数额较大”，就定罪\par}
\typeout{TCELL:p411_1:4:\the\ht0:\the\dp0}
\setbox0=\vbox{\hsize=8.440cm\parindent=0pt\linespread{1}\fontsize{10}{13.50}\selectfont\raggedright\noindent 甲向局长乙送一盒茶叶（藏有3万元现金），乙以为只是一盒普通茶叶便收下。有意见认为，乙构成受贿罪。然而，乙主观上没有收受3万元现金的故意，不构成受贿罪\par}
\typeout{TCELL:p411_1:5:\the\ht0:\the\dp0}
\setbox0=\vbox{\hsize=3.940cm\parindent=0pt\linespread{1}\fontsize{10}{13.50}\selectfont\raggedright\noindent （2）只要客观上达到“数额（特别）巨大”，就适用“数额（特别）巨大”的法定刑\par}
\typeout{TCELL:p411_1:6:\the\ht0:\the\dp0}
\setbox0=\vbox{\hsize=8.440cm\parindent=0pt\linespread{1}\fontsize{10}{13.50}\selectfont\raggedright\noindent 甲向局长乙送了一张银行卡，乙问：“卡中有多少钱？”甲担心乙拒收，便答：“3万元”（实际有20万元）。乙收下。有意见认为，对乙应适用“数额巨大”的法定刑。然而，乙主观上只有收受3万元的故意，没有收受20万元的故意，只能适用“数额较大”的法定刑，不能适用“数额巨大”的法定刑\par}
\typeout{TCELL:p411_1:7:\the\ht0:\the\dp0}
\setbox0=\vbox{\hsize=2.840cm\parindent=0pt\linespread{1}\fontsize{10}{13.50}\selectfont\raggedright\noindent 误解二：只要客观上不具备数额，就不定罪量刑\par}
\typeout{TCELL:p411_1:8:\the\ht0:\the\dp0}
\setbox0=\vbox{\hsize=3.940cm\parindent=0pt\linespread{1}\fontsize{10}{13.50}\selectfont\raggedright\noindent （1）只要客观上未达到“数额较大”，就不定罪\par}
\typeout{TCELL:p411_1:9:\the\ht0:\the\dp0}
\setbox0=\vbox{\hsize=8.440cm\parindent=0pt\linespread{1}\fontsize{10}{13.50}\selectfont\raggedright\noindent 甲真心答应向局长乙送3万元现金，乙办成事后，甲又只想送一千元，乙发现只有1千元。有意见认为，乙只收到1千元，不构成受贿罪。然而，甲开始有送3万元的意思，乙有收受3万元的故意，乙便构成受贿罪，没有收到3万元，构成受贿罪未遂。甲构成行贿罪中止。\\{}注意：如果甲一开始就不想送3万元，欺骗乙“会送3万元”，只送了1千元，那么，甲不构成行贿罪，乙不构成受贿罪，因为乙没有收到数额较大财物的可能性\par}
\typeout{TCELL:p411_1:10:\the\ht0:\the\dp0}
\setbox0=\vbox{\hsize=3.940cm\parindent=0pt\linespread{1}\fontsize{10}{13.50}\selectfont\raggedright\noindent （2）只要客观上未达到“数额（特别）巨大”，就不适用“数额（特别）巨大”的法定刑\par}
\typeout{TCELL:p411_1:11:\the\ht0:\the\dp0}
\setbox0=\vbox{\hsize=8.440cm\parindent=0pt\linespread{1}\fontsize{10}{13.50}\selectfont\raggedright\noindent 甲真心答应向局长乙送100万元现金，后甲又后悔，只送了3万元。有意见认为，对乙只适用“数额较大”的法定刑。然而，乙应构成收受“数额巨大财物”的未遂，同时构成收受“数额较大财物”的既遂，想象竞合，择一重刑论处\par}
\typeout{TCELL:p411_1:12:\the\ht0:\the\dp0}
\mbox{}
\end{document}
